% !TeX root = Linear Algebra.tex

\title{\textbf{Linear Algebra}}
\author{}
\date{}
\maketitle

\begingroup
\let\clearpage\relax
\tableofcontents
\endgroup

\clearpage

\section{Vector Spaces}\label{sec:VectorSpaces}

Linear algebra is built around linear combinations. We begin with the algebraic structures needed to define vector spaces abstractly.

\subsection{Groups and Fields}

\begin{Definition}
An \newterm{Abelian group} is an ordered pair $(G,*)$, where $G$ is a set and $*\colon G\times G\to G$ is a binary operation satisfying:
\begin{enumerate}
\item $a*(b*c)=(a*b)*c$ for all $a,b,c\in G$;
\item there exists $e\in G$ such that $a*e=e*a=a$ for all $a\in G$;
\item for every $a\in G$ there exists $b\in G$ such that $a*b=b*a=e$;
\item $a*b=b*a$ for all $a,b\in G$.
\end{enumerate}
The first three conditions define a group; the fourth is commutativity.
\end{Definition}

\begin{Example}
The sets $\mathbb R$, $\mathbb Z$, and $\mathbb Q$ form Abelian groups under addition. The sets $\mathbb Q\setminus\{0\}$ and $\mathbb R\setminus\{0\}$ form Abelian groups under multiplication. The non-zero integers do not form a group under multiplication.

For $n\in\mathbb N$, let
\begin{equation*}
\mathbb Z_n=\{0,1,\ldots,n-1\},
\end{equation*}
with addition modulo $n$. Then $(\mathbb Z_n,+_n)$ is an Abelian group. Equivalently, if $\omega=e^{2\pi i/n}$, then
\begin{equation*}
\bigl(\{1,\omega,\ldots,\omega^{n-1}\},\times\bigr)
\cong (\mathbb Z_n,+_n).
\end{equation*}
\end{Example}

\begin{Definition}
A \newterm{field} is a triple $(F,+,\cdot)$ such that:
\begin{enumerate}
\item $(F,+)$ is an Abelian group with identity $0$;
\item $(F\setminus\{0\},\cdot)$ is an Abelian group with identity $1$;
\item multiplication distributes over addition:
\begin{equation*}
a\cdot(b+c)=a\cdot b+a\cdot c
\end{equation*}
for all $a,b,c\in F$.
\end{enumerate}
\end{Definition}

The familiar sets $\mathbb Q$, $\mathbb R$, and $\mathbb C$ are fields. If $p$ is prime, then $(\mathbb Z_p,+_p,\times_p)$ is also a field. The integers with their usual addition and multiplication are a ring, but not a field.

\subsection{Definition and Basic Properties}

\begin{Definition}\label{def:VectorSpace}
Let $F$ be a field. A \newterm{vector space over $F$} is a triple $(V,+,\cdot)$, where $+\colon V\times V\to V$ is vector addition and $\cdot\colon F\times V\to V$ is scalar multiplication, satisfying:
\begin{enumerate}
\item $(V,+)$ is an Abelian group, with identity denoted by $0_V$;
\item for all $\alpha,\beta\in F$ and $u,v\in V$,
\begin{align*}
\alpha(\beta v)&=(\alpha\beta)v,\\
\alpha(u+v)&=\alpha u+\alpha v,\\
(\alpha+\beta)v&=\alpha v+\beta v;
\end{align*}
\item $1_Fv=v$ for every $v\in V$.
\end{enumerate}
The elements of $V$ are called \newterm{vectors}, and the elements of $F$ are called \newterm{scalars}.
\end{Definition}

\begin{Theorem}\label{thm:ZeroScalar}
Let $V$ be a vector space over $F$. Then, for every $v\in V$ and $\alpha\in F$,
\begin{equation*}
0_Fv=0_V,\qquad (-1_F)v=-v,\qquad \alpha 0_V=0_V.
\end{equation*}
Moreover, if $\alpha v=0_V$ and $\alpha\ne0_F$, then $v=0_V$.
\end{Theorem}
\begin{proof}
Using distributivity and the identity law in $(V,+)$,
\begin{align*}
0_Fv+0_Fv&=(0_F+0_F)v=0_Fv,\\
v+(-1_F)v&=(1_F-1_F)v=0_Fv=0_V.
\end{align*}
Cancellation gives $0_Fv=0_V$ and $(-1_F)v=-v$. Similarly,
\begin{equation*}
\alpha0_V=\alpha(0_V+0_V)=\alpha0_V+\alpha0_V,
\end{equation*}
so $\alpha0_V=0_V$. Finally, if $\alpha\ne0_F$ and $\alpha v=0_V$, then
\begin{equation*}
v=1_Fv=\alpha^{-1}(\alpha v)=\alpha^{-1}0_V=0_V.
\end{equation*}
\end{proof}

\begin{Corollary}
For $\alpha,\beta\in F$ and $v\in V$,
\begin{equation*}
(\alpha-\beta)v=\alpha v-\beta v.
\end{equation*}
\end{Corollary}

\begin{Example}\label{ex:TupleSpace}
For a field $F$ and $n\in\mathbb N$, define
\begin{equation*}
F^n=\qty{\, (x_1,\ldots,x_n) \mid x_i\in F\text{ for }i=1,\ldots,n \,}.
\end{equation*}
For $x=(x_1,\ldots,x_n)$ and $y=(y_1,\ldots,y_n)$, define
\begin{align*}
x+y&=(x_1+y_1,\ldots,x_n+y_n),\\
\alpha x&=(\alpha x_1,\ldots,\alpha x_n).
\end{align*}
Then $F^n$ is a vector space over $F$. Its zero vector is $(0,\ldots,0)$.
\end{Example}

\begin{Example}
Taking $F=\mathbb Z_2=\{0,1\}$, the vector space $F^3$ has eight elements:
\begin{equation*}
F^3=\{(0,0,0),(0,0,1),\ldots,(1,1,1)\}.
\end{equation*}
It may be visualised as the vertices of a cube, with vector addition performed coordinatewise modulo $2$.
\end{Example}

\begin{Example}
The following are vector spaces with their indicated scalar fields:
\begin{enumerate}
\item $\mathbb C^n$ over $\mathbb C$;
\item $\mathbb C^n$ over $\mathbb R$, using the usual real scalar multiplication;
\item $\mathbb R^n$ over $\mathbb R$.
\end{enumerate}
With the usual scalar multiplication, $\mathbb R^n$ is not a vector space over $\mathbb C$. However, $\mathbb R^2$ can be made into a vector space over $\mathbb C$ by defining, for $z=u+iv\in\mathbb C$,
\begin{equation*}
z(x_1,x_2)=(ux_1-vx_2,\;vx_1+ux_2).
\end{equation*}
This is the usual identification of $\mathbb R^2$ with $\mathbb C$.
\end{Example}

\begin{Example}\label{ex:MatrixSpace}
Let $\mathbb F^{m\times n}$ be the set of all $m\times n$ matrices over $\mathbb F$. For $A=[a_{ij}]$ and $B=[b_{ij}]$, define
\begin{equation*}
A+B=[a_{ij}+b_{ij}],\qquad \alpha A=[\alpha a_{ij}].
\end{equation*}
Then $\mathbb F^{m\times n}$ is a vector space over $\mathbb F$.
\end{Example}

\begin{Example}\label{ex:FunctionSpace}
Let $S$ be a non-empty set and $F$ a field. The set
\begin{equation*}
F^S=\qty{\, f \mid f\colon S\to F \,}
\end{equation*}
is a vector space over $F$ under pointwise operations:
\begin{align*}
(f+g)(x)&=f(x)+g(x),\\
(\alpha f)(x)&=\alpha f(x),
\end{align*}
for every $x\in S$. When $S=\{1,\ldots,n\}$, this is essentially the same vector space as $F^n$.
\end{Example}

\begin{Example}
The set of continuous functions $\mathbb R\to\mathbb R$, and the set of differentiable functions $\mathbb R\to\mathbb R$, are vector spaces over $\mathbb R$. More generally, the solution set of a homogeneous linear differential equation
\begin{equation*}
\frac{d^ny}{dx^n}+a_1\frac{d^{n-1}y}{dx^{n-1}}+\cdots+a_ny=0,
\qquad a_1,\ldots,a_n\in\mathbb R,
\end{equation*}
is a vector space over $\mathbb R$.
\end{Example}

\subsection{Linear Combinations}

\begin{Definition}
Let $V$ be a vector space over $F$ and let $v_1,\ldots,v_k\in V$. A vector
\begin{equation*}
\alpha_1v_1+\cdots+\alpha_kv_k,
\qquad \alpha_1,\ldots,\alpha_k\in F,
\end{equation*}
is a \newterm{linear combination} of $v_1,\ldots,v_k$.
\end{Definition}

\section{Subspaces}\label{sec:Subspaces}

\begin{Definition}
Let $V$ be a vector space over $F$. A subset $W\subseteq V$ is a \newterm{subspace} of $V$ if $W$ is itself a vector space over $F$ under the operations inherited from $V$.
\end{Definition}

\begin{Example}
In $\mathbb R^2$, the sets
\begin{equation*}
X=\qty{\, \begin{bmatrix}x\\0\end{bmatrix} \mid x\in\mathbb R \,},\qquad
Y=\qty{\, \begin{bmatrix}0\\y\end{bmatrix} \mid y\in\mathbb R \,},\qquad
W=\qty{\, \begin{bmatrix}x\\mx\end{bmatrix} \mid x\in\mathbb R \,}
\end{equation*}
are subspaces for every $m\in\mathbb R$. For example,
\begin{align*}
(x_1,mx_1)+(x_2,mx_2)&=(x_1+x_2,m(x_1+x_2)),\\
\alpha(x,mx)&=(\alpha x,m(\alpha x)).
\end{align*}
\end{Example}

The set $\{0\}$ is the \newterm{zero subspace} (or trivial subspace) of any vector space. The whole vector space $V$ is also a subspace. In contrast,
\begin{equation*}
S=\qty{\, \begin{bmatrix}x\\1\end{bmatrix} \mid x\in\mathbb R \,}\subseteq\mathbb R^2
\end{equation*}
is not a subspace, since it does not contain the zero vector.

\begin{Theorem}\label{thm:SubspaceCriterion}
A subset $W$ of a vector space $V$ over $F$ is a subspace if and only if
\begin{equation*}
W\ne\varnothing\quad\text{and}\quad
\alpha u+v\in W\quad\text{for all }\alpha\in F,\ u,v\in W.
\end{equation*}
\end{Theorem}
\begin{proof}
If $W$ is a subspace, the condition is immediate. Conversely, suppose that $W\ne\varnothing$ and that the displayed closure condition holds. Choose $w\in W$. Then $(-1)w+w=0_V\in W$. Hence, for $u,v\in W$, we have $u+v=1u+v\in W$; for $\alpha\in F$ and $u\in W$, we have $\alpha u+0_V\in W$; and $-u=(-1)u\in W$. Thus $W$ is closed under the inherited operations and contains the additive identity and inverses. All remaining vector-space axioms are inherited from $V$.
\end{proof}

\begin{Theorem}\label{thm:Intersection}
Let $\{W_i\}_{i\in I}$ be a collection of subspaces of $V$. Then
\begin{equation*}
\bigcap_{i\in I}W_i
\end{equation*}
is also a subspace of $V$.
\end{Theorem}
\begin{proof}
The zero vector belongs to every $W_i$, so the intersection is non-empty. If $u,v$ belong to the intersection and $\alpha\in F$, then $u,v\in W_i$ for every $i\in I$. By the subspace criterion, $\alpha u+v\in W_i$ for every $i\in I$, and hence $\alpha u+v$ belongs to the intersection.
\end{proof}

\begin{Theorem}
If $W_1$ and $W_2$ are subspaces of $V$, then $W_1\cup W_2$ is a subspace if and only if $W_1\subseteq W_2$ or $W_2\subseteq W_1$.
\end{Theorem}
\begin{proof}
If one subspace contains the other, the union is the larger subspace. Conversely, assume neither contains the other. Choose $w_1\in W_1\setminus W_2$ and $w_2\in W_2\setminus W_1$. If $w_1+w_2\in W_1$, then $w_2=(w_1+w_2)-w_1\in W_1$, a contradiction. Similarly, $w_1+w_2\notin W_2$. Thus the union is not closed under addition.
\end{proof}

\begin{Definition}\label{def:Span}
Let $S=\{v_i\}_{i\in I}\subseteq V$. The \newterm{span} of $S$ is the set of all finite linear combinations of elements of $S$:
\begin{align*}
\operatorname{span}S
&=\qty{\, \sum_{i\in I}\alpha_iv_i \mid
\alpha_i\in F\text{ and all but finitely many }\alpha_i=0 \,}.
\end{align*}
Equivalently,
\begin{equation*}
\operatorname{span}S
=\bigcap\qty{\, W\subseteq V \mid S\subseteq W\text{ and }W\text{ is a subspace of }V \,}.
\end{equation*}
Thus $\operatorname{span}S$ is the smallest subspace of $V$ containing $S$. The set $S$ is a \newterm{spanning set} of $V$ if $\operatorname{span}S=V$.
\end{Definition}

\begin{Example}
Let $V=\mathbb R^3$ and consider
\begin{align*}
S_1&=\left\{\begin{bmatrix}1\\-1\\0\end{bmatrix},
\begin{bmatrix}3\\-3\\0\end{bmatrix},
\begin{bmatrix}-1\\1\\0\end{bmatrix}\right\},\\
S_2&=\left\{\begin{bmatrix}1\\0\\0\end{bmatrix},
\begin{bmatrix}0\\1\\0\end{bmatrix},
\begin{bmatrix}0\\0\\1\end{bmatrix}\right\},\\
S_3&=\left\{\begin{bmatrix}1\\1\\1\end{bmatrix},
\begin{bmatrix}1\\-1\\2\end{bmatrix},
\begin{bmatrix}1\\3\\0\end{bmatrix}\right\}.
\end{align*}
Then
\begin{align*}
\operatorname{span}S_1
 &=\qty{\, \begin{bmatrix}x\\-x\\0\end{bmatrix} \mid x\in\mathbb R\,},\\
\operatorname{span}S_2&=\mathbb R^3,\\
\operatorname{span}S_3
 &=\qty{\, \begin{bmatrix}x+y\\x-y\\x+2y\end{bmatrix}
      \mid x,y\in\mathbb R\,}\ne\mathbb R^3.
\end{align*}
Indeed, $\begin{bmatrix}1\\3\\0\end{bmatrix}= 2\begin{bmatrix}1\\1\\1\end{bmatrix}-\begin{bmatrix}1\\-1\\2\end{bmatrix}$.
\end{Example}

\begin{Example}
Let $V$ be the vector space of all real polynomials and let
\begin{equation*}
S=\{1,x,x^2,\ldots\}.
\end{equation*}
Then $\operatorname{span}S=\mathbb R[x]$, the vector space of polynomials. This is not the vector space of all real-valued functions of $x$.
\end{Example}

\section{Linear Independence and Bases}\label{sec:Bases}

\begin{Definition}
Let $S=\{v_i\}_{i\in I}\subseteq V$. The set $S$ is \newterm{linearly independent} if, for every finite selection $v_{i_1},\ldots,v_{i_k}$ of distinct elements of $S$,
\begin{equation*}
\alpha_1v_{i_1}+\cdots+\alpha_kv_{i_k}=0_V
\quad\Longrightarrow\quad
\alpha_1=\cdots=\alpha_k=0_F.
\end{equation*}
If $S$ is not linearly independent, it is \newterm{linearly dependent}.
\end{Definition}

\begin{Definition}
A \newterm{basis} of a vector space $V$ is a linearly independent spanning set of $V$.
\end{Definition}

\begin{Example}
The standard basis of $\mathbb R^3$ is
\begin{equation*}
\mathcal E=\left\{\begin{bmatrix}1\\0\\0\end{bmatrix},
\begin{bmatrix}0\\1\\0\end{bmatrix},
\begin{bmatrix}0\\0\\1\end{bmatrix}\right\}.
\end{equation*}
The set of the first two vectors is linearly independent but does not span $\mathbb R^3$. A set of four vectors in $\mathbb R^3$ cannot be linearly independent. The set
\begin{equation*}
\left\{\begin{bmatrix}1\\1\end{bmatrix},
\begin{bmatrix}1\\-1\end{bmatrix}\right\}
\end{equation*}
is a basis of $\mathbb R^2$, and $\{1,x,x^2,\ldots\}$ is a basis of $\mathbb R[x]$.
\end{Example}

\begin{Example}
In $\mathbb R^3$, let $v_1,v_2,v_3$ be the standard basis vectors and put
\begin{equation*}
v_4=\begin{bmatrix}1\\0\\2\end{bmatrix},\qquad
S_1=\{v_1,v_2,v_3,v_4\}.
\end{equation*}
Since $v_1+2v_3-v_4=0$, the set $S_1$ is dependent. Now let
\begin{equation*}
S_2=\left\{\begin{bmatrix}1\\0\\0\end{bmatrix},
\begin{bmatrix}0\\1\\0\end{bmatrix}\right\}.
\end{equation*}
This set is independent but does not span $\mathbb R^3$. However, $S_2\cup\{(0,0,1)\}$ is a basis of $\mathbb R^3$.
\end{Example}

\begin{Theorem}\label{thm:DependentVector}
Let $S=\{v_1,\ldots,v_n\}\subseteq V$. Then $S$ is linearly dependent if and only if some vector of $S$ can be written as a linear combination of the other vectors.
\end{Theorem}
\begin{proof}
If, say, $v_n=\alpha_1v_1+\cdots+\alpha_{n-1}v_{n-1}$, then
\begin{equation*}
\alpha_1v_1+\cdots+\alpha_{n-1}v_{n-1}-v_n=0_V
\end{equation*}
is a non-trivial linear relation. Thus $S$ is dependent.

Conversely, suppose that $S$ is dependent. Then there are scalars $\alpha_1,\ldots,\alpha_n$, not all zero, such that
\begin{equation*}
\alpha_1v_1+\cdots+\alpha_nv_n=0_V.
\end{equation*}
Choose $i$ with $\alpha_i\ne0_F$. Rearranging and dividing by $\alpha_i$ gives
\begin{equation*}
v_i=-\frac{\alpha_1}{\alpha_i}v_1-\cdots-
\frac{\alpha_{i-1}}{\alpha_i}v_{i-1}
-\frac{\alpha_{i+1}}{\alpha_i}v_{i+1}-\cdots-
\frac{\alpha_n}{\alpha_i}v_n.
\end{equation*}
\end{proof}

\subsection{Finite Dimensional Vector Spaces}

\begin{Definition}
A vector space $V$ is \newterm{finite dimensional} if it has a finite basis. If it has no finite basis, it is \newterm{infinite dimensional}.
\end{Definition}

\begin{Example}
The vector space $F^n$ has the standard basis
\begin{equation*}
e_i=(0,\ldots,0,1,0,\ldots,0),
\end{equation*}
where the $1$ is in the $i$\nth component. Hence $F^n$ is finite dimensional.

The polynomial space $F[x]$ is infinite dimensional. If $p_1,\ldots,p_n\in F[x]$ and $k$ is the largest of their degrees, then every linear combination of these polynomials has degree at most $k$, whereas $x^{k+1}$ does not. Therefore $\{1,x,x^2,\ldots\}$ has no finite spanning subset.
\end{Example}

\begin{Theorem}\label{thm:FiniteSpanningBasis}
If $V$ is spanned by a finite set $S$, then $V$ has a basis contained in $S$.
\end{Theorem}
\begin{proof}
If $S$ is independent, it is already a basis. Otherwise, by \cref{thm:DependentVector}, one element of $S$ is a linear combination of the others. Remove that element. The remaining set still spans $V$. Repeating this process must terminate because $S$ is finite, and it terminates at a linearly independent spanning set.
\end{proof}

\begin{Theorem}\label{thm:BasisCardinality}
Any two bases of a finite dimensional vector space have the same number of elements.
\end{Theorem}
\begin{proof}
We first prove the replacement principle: if $u_1,\ldots,u_r$ are linearly independent and $v_1,\ldots,v_n$ span $V$, then $r\le n$. Express $u_1$ in terms of the spanning vectors. Some coefficient is non-zero, so we can solve for the corresponding $v_j$ and replace it by $u_1$ without changing the span. After $k$ replacements we have a spanning list consisting of $u_1,\ldots,u_k$ and $n-k$ of the original vectors. In an expression for $u_{k+1}$ using this list, some coefficient of an original vector must be non-zero; otherwise $u_{k+1}$ belongs to $\operatorname{span}\{u_1,\ldots,u_k\}$, contrary to independence. Thus another replacement is possible. There cannot be more than $n$ replacements, giving $r\le n$.

If $V$ has a basis of size $n$, no independent set can have more than $n$ elements, since any $n+1$ of them would contradict the replacement principle. In particular every basis is finite. Applying the principle in both directions to bases of sizes $m$ and $n$ gives $m\le n$ and $n\le m$, so $m=n$.
\end{proof}

\begin{Theorem}\label{thm:BasisExtension}
Every linearly independent subset of a finite dimensional vector space can be extended to a basis.
\end{Theorem}
\begin{proof}
Let $S$ be independent, and let $v_1,\ldots,v_n$ be a finite spanning list. Consider these vectors in turn, adjoining $v_j$ to $S$ precisely when it is not in the span of the vectors already selected. Independence is preserved at every step. The resulting set spans every $v_j$, hence spans $V$, and is therefore a basis.
\end{proof}

\subsection{Dimension and Coordinates}

\begin{Definition}
If $V$ is finite dimensional, its \newterm{dimension}, denoted by $\dim V$, is the number of vectors in any basis of $V$.
\end{Definition}

\begin{Example}
\begin{enumerate}
\item $\dim F^n=n$.
\item If $W=\qty{\, \begin{bmatrix}x\\0\\0\end{bmatrix} \mid x\in\mathbb R \,}\subseteq\mathbb R^3$, then $\dim W=1$.
\item If $U=\qty{\, \begin{bmatrix}x\\y\\0\end{bmatrix} \mid x,y\in\mathbb R \,}\subseteq\mathbb R^3$, then $\dim U=2$.
\item $\dim\{0_V\}=0$.
\item Every field $F$ is a one-dimensional vector space over itself, with basis $\{1_F\}$.
\end{enumerate}
\end{Example}

\begin{Theorem}\label{thm:Coordinates}
Let $\mathcal B=\{v_i\}_{i\in I}$ be a basis of $V$. Every vector of $V$ can be written as a linear combination of elements of $\mathcal B$ in exactly one way.
\end{Theorem}
\begin{proof}
Existence follows from the fact that $\mathcal B$ spans $V$. Suppose
\begin{equation*}
v=\sum_{k=1}^n\alpha_kv_{i_k}
 =\sum_{\ell=1}^m\beta_\ell v_{j_\ell}
\end{equation*}
are two representations. Moving all terms to one side and collecting equal basis vectors gives a finite linear combination of elements of $\mathcal B$ equal to $0_V$. Linear independence implies that every coefficient is zero, so the representations are identical.
\end{proof}

If $\mathcal B=(v_1,\ldots,v_n)$ is an ordered basis and $v=\alpha_1v_1+\cdots+\alpha_nv_n$, then
\begin{equation*}
[v]_{\mathcal B}=\begin{bmatrix}\alpha_1\\\vdots\\\alpha_n\end{bmatrix}
\end{equation*}
is the \newterm{coordinate representation} of $v$ with respect to $\mathcal B$.

\begin{Example}
Let $V=\mathbb R^3$ and take
\begin{align*}
\mathcal B_1&=\left\{\begin{bmatrix}1\\0\\0\end{bmatrix},
\begin{bmatrix}0\\1\\0\end{bmatrix},
\begin{bmatrix}0\\0\\1\end{bmatrix}\right\},\\
\mathcal B_2&=\left\{\begin{bmatrix}1\\1\\0\end{bmatrix},
\begin{bmatrix}0\\1\\0\end{bmatrix},
\begin{bmatrix}0\\0\\1\end{bmatrix}\right\}.
\end{align*}
For $a=(x,y,z)$,
\begin{equation*}
[a]_{\mathcal B_1}=(x,y,z),\qquad
[a]_{\mathcal B_2}=(x,y-x,z).
\end{equation*}
\end{Example}

\begin{Example}
For $V=F[x]$ with ordered basis $\mathcal B=(1,x,x^2,\ldots)$, if
\begin{equation*}
p(x)=a_0+a_1x+\cdots+a_nx^n,
\end{equation*}
then
\begin{equation*}
[p]_{\mathcal B}=(a_0,a_1,\ldots,a_n,0,0,\ldots).
\end{equation*}
\end{Example}

% !TeX root = Linear Algebra.tex
% Supplementary material from handwritten pages 1--8.

\begin{Corollary}\label{cor:DimensionTests}
Let $\dim V=n<\infty$ and let $W\subseteq V$ be a subspace. Then $W$ is finite dimensional, $\dim W\le n$, and equality holds if and only if $W=V$. A set of more than $n$ vectors is dependent, and a set of fewer than $n$ vectors cannot span $V$. For a set of exactly $n$ vectors, being a basis, being linearly independent, and spanning $V$ are equivalent.
\end{Corollary}
\begin{proof}
Choose vectors of $W$ successively outside the span of those already selected. The replacement principle ensures that this process stops after at most $n$ steps, producing a basis of $W$. Extend it to a basis of $V$ using \cref{thm:BasisExtension}. The assertions follow from basis extension, extraction of a basis from a finite spanning set, and \cref{thm:BasisCardinality}.
\end{proof}

\begin{Example}
Write $P_n(F)$ for the space of polynomials over $F$ of degree at most $n$, including the zero polynomial. Its standard basis is $(1,x,\ldots,x^n)$, so $\dim P_n(F)=n+1$. The matrices $E_{ij}$, with a $1$ in position $(i,j)$ and zeros elsewhere, form the standard basis of $F^{m\times n}$, since
\begin{equation*}
A=\sum_{i=1}^m\sum_{j=1}^n a_{ij}E_{ij}
\end{equation*}
uniquely. Thus $\dim F^{m\times n}=mn$. Also, $\{1,i\}$ is a basis of $\mathbb C$ over $\mathbb R$, whereas $\{1\}$ is a basis over $\mathbb C$. Dimension therefore depends on the scalar field.
\end{Example}

\begin{Definition}
An \newterm{ordered basis} is a basis with a specified order. If $\mathcal B=(v_1,\ldots,v_n)$ and $v=\sum_{i=1}^n\alpha_iv_i$, then $\alpha_i$ is the $i$\nth \newterm{coordinate} (or \newterm{component}) of $v$ relative to $\mathcal B$.
\end{Definition}

\begin{Example}
For $v=\begin{bmatrix}2\\-3\end{bmatrix}$ in $\mathbb R^2$, the coordinates in the standard basis are $[v]_{\mathcal E}=\begin{bmatrix}2\\-3\end{bmatrix}$. In the ordered basis $\mathcal B=\bigl(\begin{bmatrix}1\\1\end{bmatrix},\begin{bmatrix}1\\-1\end{bmatrix}\bigr)$, they are
\begin{equation*}
[v]_{\mathcal B}=\begin{bmatrix}-1/2\\5/2\end{bmatrix}.
\end{equation*}
For the solution space of $y''-y=0$, two ordered bases are $\mathcal B_1=(e^x,e^{-x})$ and $\mathcal B_2=(\cosh x,\sinh x)$. Since $3e^x-5e^{-x}=-2\cosh x+8\sinh x$, its coordinate vectors are $\begin{bmatrix}3\\-5\end{bmatrix}$ and $\begin{bmatrix}-2\\8\end{bmatrix}$ respectively.
\end{Example}

\subsection{Further Examples and Exercises}

\begin{Exercise}
With usual addition, show that neither of the following defines a real vector space: $\mathbb R^n$ for $n\ge2$ with $\alpha\cdot(x_1,\ldots,x_n)=(\alpha x_1,\alpha^2x_2,\ldots,\alpha^nx_n)$; and $\mathbb R^2$ with $\alpha\cdot(x,y)=(\alpha x,0)$. Identify an axiom that fails in each case.
\end{Exercise}

\begin{Exercise}
Show that $\qty{\, A=[a_{ij}]\in F^{m\times n} \mid a_{pq}=0 \,}$ is a subspace for any fixed position $(p,q)$. Explain why $\mathbb R^2$ is a real subspace of $\mathbb C^2$, but not a complex subspace. Find all subspaces of $\mathbb R^2$.
\end{Exercise}

\begin{Exercise}
Prove directly that $\operatorname{span}S$ is a subspace and is contained in every subspace containing $S$. Use the convention that the empty linear combination is $0$ to show $\operatorname{span}\varnothing=\{0\}$. Prove that $S_1\subseteq S_2$ implies $\operatorname{span}S_1\subseteq\operatorname{span}S_2$.
\end{Exercise}

\begin{Example}
In $\mathbb R^2$, each of the sets
\begin{align*}
S_1&=\qty{\begin{bmatrix}1\\0\end{bmatrix},\begin{bmatrix}0\\1\end{bmatrix}},&
S_2&=\qty{\begin{bmatrix}1\\0\end{bmatrix},\begin{bmatrix}0\\-1\end{bmatrix}},\\
S_3&=\qty{\begin{bmatrix}1\\1\end{bmatrix},\begin{bmatrix}1\\-1\end{bmatrix}},&
S_4&=\qty{\begin{bmatrix}1\\0\end{bmatrix},\begin{bmatrix}0\\1\end{bmatrix},\begin{bmatrix}1\\1\end{bmatrix}}
\end{align*}
spans $\mathbb R^2$, but $S_4$ is dependent. In contrast,
\begin{align*}
\operatorname{span}\qty{\begin{bmatrix}1\\0\end{bmatrix},\begin{bmatrix}-1\\0\end{bmatrix}}&=\qty{\, \begin{bmatrix}x\\0\end{bmatrix} \mid x\in\mathbb R \,},\\
\operatorname{span}\qty{\begin{bmatrix}2\\4\end{bmatrix},\begin{bmatrix}3\\6\end{bmatrix}}&=\qty{\, \begin{bmatrix}x\\2x\end{bmatrix} \mid x\in\mathbb R \,}.
\end{align*}
The empty set is a basis of $\{0\}$. A set containing $0$ is always linearly dependent.
\end{Example}

\begin{Exercise}
Prove the converse of \cref{thm:Coordinates}: if every vector has a unique expression as a finite linear combination of elements of $\mathcal B$, then $\mathcal B$ is a basis. Here uniqueness concerns the coefficients of each basis vector; adding terms with coefficient zero does not give a different expression.
\end{Exercise}

\begin{Exercise}
Let $V$ be the space of functions $f\colon\mathbb R\to\mathbb R$ which vanish at all but finitely many points. For each $a\in\mathbb R$, define
\begin{equation*}
\delta_a(x)=\begin{cases}1,&x=a,\\0,&x\ne a.\end{cases}
\end{equation*}
Show that $V$ is a vector space under pointwise operations and that $\qty{\, \delta_a \mid a\in\mathbb R \,}$ is a basis. Notice that every individual function has finite support even though the basis is uncountable.
\end{Exercise}


\section{Linear Transformations}\label{sec:LinearTransformations}

\begin{Definition}
Let $V$ and $W$ be vector spaces over the same field $F$. A function $T\colon V\to W$ is a \newterm{linear transformation} if
\begin{equation*}
T(\alpha u+v)=\alpha T(u)+T(v)
\end{equation*}
for all $\alpha\in F$ and $u,v\in V$.
\end{Definition}

By induction, a linear transformation preserves all finite linear combinations:
\begin{equation*}
T(\alpha_1v_1+\cdots+\alpha_nv_n)
=\alpha_1T(v_1)+\cdots+\alpha_nT(v_n).
\end{equation*}
It also satisfies $T(0_V)=0_W$, and is a homomorphism from $(V,+)$ to $(W,+)$.

The \newterm{image} (or range) of $T$ is
\begin{equation*}
\im T=T(V)=\qty{\, T(v) \mid v\in V \,}\subseteq W,
\end{equation*}
and the \newterm{kernel} (or null space) of $T$ is
\begin{equation*}
\ker T=\qty{\, v\in V \mid T(v)=0_W \,}\subseteq V.
\end{equation*}

\begin{Theorem}\label{thm:ImageKernel}
If $T\colon V\to W$ is linear, then $\im T$ is a subspace of $W$ and $\ker T$ is a subspace of $V$.
\end{Theorem}
\begin{proof}
Since $T(0_V)=0_W$, the image is non-empty. If $x=T(u)$ and $y=T(v)$ belong to $\im T$, then
\begin{equation*}
\alpha x+y=\alpha T(u)+T(v)=T(\alpha u+v)\in\im T.
\end{equation*}
Thus $\im T$ is a subspace by the subspace criterion.

Also $0_V\in\ker T$. If $u,v\in\ker T$, then
\begin{equation*}
T(\alpha u+v)=\alpha T(u)+T(v)=0_W,
\end{equation*}
so $\alpha u+v\in\ker T$. Hence $\ker T$ is a subspace.
\end{proof}

\begin{Example}
Define $T\colon\mathbb R^3\to\mathbb R^2$ by
\begin{equation*}
T(x,y,z)=(z,x+y).
\end{equation*}
Then
\begin{align*}
\ker T
 &=\qty{\, (x,y,z)\in\mathbb R^3 \mid (z,x+y)=(0,0) \,}\\
 &=\qty{\, \begin{bmatrix}x\\-x\\0\end{bmatrix} \mid x\in\mathbb R \,},\\
\im T
 &=\qty{\, (z,x+y) \mid x,y,z\in\mathbb R\,}=\mathbb R^2.
\end{align*}
Therefore $\dim\ker T=1$, $\dim\im T=2$, and $\dim\mathbb R^3=3$.
\end{Example}

\begin{Example}
Define $T\colon\mathbb R^3\to\mathbb R^2$ by
\begin{equation*}
T(x,y,z)=(y,-y).
\end{equation*}
Then
\begin{align*}
\ker T&=\qty{\, \begin{bmatrix}x\\0\\z\end{bmatrix} \mid x,z\in\mathbb R \,},\\
\im T&=\qty{\, \begin{bmatrix}a\\-a\end{bmatrix} \mid a\in\mathbb R \,}.
\end{align*}
Thus $\dim\ker T=2$, $\dim\im T=1$, and $\dim\mathbb R^3=3$.
\end{Example}

\begin{Example}
Define $T\colon\mathbb R^2\to\mathbb R^3$ by
\begin{equation*}
T(x,y)=(x+y,x,y+x).
\end{equation*}
Then
\begin{align*}
\ker T&=\qty{\begin{bmatrix}0\\0\end{bmatrix}},\\
\im T&=\qty{\, \begin{bmatrix}x+y\\x\\x+y\end{bmatrix} \mid x,y\in\mathbb R \,}\\
 &=\qty{\, \begin{bmatrix}a\\b\\a\end{bmatrix} \mid a,b\in\mathbb R \,}\\
 &=\operatorname{span}\qty{\begin{bmatrix}1\\0\\1\end{bmatrix},\begin{bmatrix}0\\1\\0\end{bmatrix}}.
\end{align*}
Hence $\dim\ker T=0$, $\dim\im T=2$, and $\dim\mathbb R^2=2$.
\end{Example}

\begin{Example}
Let $V$ be the vector space of all infinitely differentiable functions $\mathbb R\to\mathbb R$. Define the differentiation operator
\begin{equation*}
D\colon V\to V,\qquad Df=f'.
\end{equation*}
For $\alpha\in\mathbb R$ and $f,g\in V$,
\begin{equation*}
D(\alpha f+g)=(\alpha f+g)'=\alpha f'+g'=\alpha Df+Dg,
\end{equation*}
so $D$ is linear. Its kernel is the space of constant functions:
\begin{equation*}
\ker D=\qty{\, f\in V \mid f'=0 \,}=\qty{\, x\mapsto c \mid c\in\mathbb R\,}.
\end{equation*}
\end{Example}

\begin{Example}
Let $V=\mathbb R[x]$ and let $D\colon V\to V$ again be differentiation. If
\begin{equation*}
p(x)=a_0+a_1x+\cdots+a_nx^n,
\end{equation*}
then
\begin{equation*}
Dp(x)=a_1+2a_2x+3a_3x^2+\cdots+na_nx^{n-1}.
\end{equation*}
Therefore
\begin{equation*}
\ker D=\qty{\, c \mid c\in\mathbb R\,},\qquad \im D=\mathbb R[x]=V,
\end{equation*}
since every polynomial has a polynomial antiderivative.
\end{Example}

% !TeX root = Linear Algebra.tex
% Handwritten pages 8--20, with repeated examples consolidated.

\subsection{Examples and Rank--Nullity}

\begin{Example}
The identity $\id_V(v)=v$, the zero transformation $0(v)=0_W$, and scalar multiplication $v\mapsto\alpha v$ are linear. Over $\mathbb R$, the maps $\mathbb C\to\mathbb R^2$, $x+iy\mapsto(x,y)$, and $\mathbb R^2\to\mathbb C$, $(x,y)\mapsto x+iy$, are inverse linear transformations.

A map $\mathbb R^2\to\mathbb R^3$ of the form $(x,y)\mapsto(a_1x+b_1y,a_2x+b_2y,a_3x+b_3y)$ is linear. The map $(x,y)\mapsto(x^2,0,0)$ is not linear, since it does not preserve scalar multiplication. The map $(x,y)\mapsto(x+1,y+1,x+y)$ is not linear, since it does not take $0$ to $0$.
\end{Example}

\begin{Definition}
The \newterm{rank} and \newterm{nullity} of a linear transformation $T$ are
\begin{equation*}
\rank T=\dim\im T,\qquad \nullity T=\dim\ker T.
\end{equation*}
\end{Definition}

\begin{Theorem}[Rank--Nullity Theorem]\label{thm:RankNullity}
If $T\colon V\to W$ is linear and $V$ is finite dimensional, then
\begin{equation*}
\rank T+\nullity T=\dim V.
\end{equation*}
\end{Theorem}
\begin{proof}
Choose a basis $v_1,\ldots,v_k$ of $\ker T$ and extend it to a basis $v_1,\ldots,v_n$ of $V$. We claim that $T(v_{k+1}),\ldots,T(v_n)$ is a basis of $\im T$. For $v=\sum_{i=1}^n\alpha_iv_i$, linearity gives
\begin{equation*}
T(v)=\sum_{i=k+1}^n\alpha_iT(v_i),
\end{equation*}
so these vectors span the image. If $\sum_{i=k+1}^n\beta_iT(v_i)=0$, then $\sum_{i=k+1}^n\beta_iv_i\in\ker T$. Expressing this vector using $v_1,\ldots,v_k$ gives a linear relation among the basis vectors of $V$. All coefficients are zero, proving independence. Consequently $\rank T=n-k$ and $\nullity T=k$.
\end{proof}

\begin{Example}
For $T\colon\mathbb R^2\to\mathbb R^2$ given by $T(x,y)=(x-y,y-x)$,
\begin{equation*}
\ker T=\operatorname{span}\qty{\begin{bmatrix}1\\1\end{bmatrix}},\qquad \im T=\operatorname{span}\qty{\begin{bmatrix}1\\-1\end{bmatrix}}.
\end{equation*}
For the following maps, the kernels and images can similarly be read off from the formulas:
\begin{align*}
T_1(x,y)&=(x,y,x-y,x+y),& \ker T_1&=\{0\},\\
\im T_1&=\operatorname{span}\qty{\begin{bmatrix}1\\0\\1\\1\end{bmatrix},\begin{bmatrix}0\\1\\-1\\1\end{bmatrix}};\\
T_2(x,y)&=(x-y,0,y-x,2x-2y),& \ker T_2&=\operatorname{span}\qty{\begin{bmatrix}1\\1\end{bmatrix}},\\
\im T_2&=\operatorname{span}\qty{\begin{bmatrix}1\\0\\-1\\2\end{bmatrix}};\\
T_3(x,y,z)&=(x+y,x-y,0,x),& \ker T_3&=\operatorname{span}\qty{\begin{bmatrix}0\\0\\1\end{bmatrix}},\\
\im T_3&=\operatorname{span}\qty{\begin{bmatrix}1\\1\\0\\1\end{bmatrix},\begin{bmatrix}1\\-1\\0\\0\end{bmatrix}}.
\end{align*}
The respective rank--nullity pairs are $(2,0)$, $(1,1)$ and $(2,1)$.
\end{Example}

\begin{Theorem}\label{thm:InjectiveKernel}
A linear transformation $T\colon V\to W$ is injective if and only if $\ker T=\{0_V\}$.
\end{Theorem}
\begin{proof}
If $T$ is injective, $T(v)=0_W=T(0_V)$ implies $v=0_V$. Conversely, if the kernel is trivial and $T(u)=T(v)$, then $T(u-v)=0_W$, so $u=v$.
\end{proof}

\begin{Exercise}
Prove that an injective linear transformation takes every linearly independent set to a linearly independent set, and a surjective linear transformation takes every spanning set to a spanning set. Deduce, in finite dimensions, that injectivity implies $\dim V\le\dim W$, while surjectivity implies $\dim V\ge\dim W$.
\end{Exercise}

\subsection{Specifying a Transformation on a Basis}

\begin{Theorem}\label{thm:SpecifyBasis}
Let $\mathcal B=(v_1,\ldots,v_n)$ be a basis of $V$, and choose arbitrary vectors $w_1,\ldots,w_n\in W$. There is exactly one linear transformation $T\colon V\to W$ such that $T(v_i)=w_i$ for every $i$.
\end{Theorem}
\begin{proof}
The unique basis expansion of each vector makes the formula
\begin{equation*}
T\pqty{\sum_{i=1}^n\alpha_iv_i}=\sum_{i=1}^n\alpha_iw_i
\end{equation*}
well defined. If $u=\sum_i\alpha_iv_i$ and $v=\sum_i\beta_iv_i$, then
\begin{equation*}
T(\gamma u+v)=\sum_i(\gamma\alpha_i+\beta_i)w_i=\gamma T(u)+T(v).
\end{equation*}
Thus $T$ is linear and has the required values. Every linear transformation with these values must satisfy the displayed defining formula, proving uniqueness.
\end{proof}

\begin{Example}
If $T\colon\mathbb R^2\to\mathbb R^2$ satisfies $T(1,0)=(-1,2)$ and $T(0,1)=(1,3)$, then $T(x,y)=(y-x,2x+3y)$. If instead $T\colon\mathbb R^2\to\mathbb R^3$ is specified by $T(1,1)=(1,0,1)$ and $T(1,-1)=(0,2,-1)$, then
\begin{equation*}
T(x,y)=\frac{x+y}{2}\begin{bmatrix}1\\0\\1\end{bmatrix}+\frac{x-y}{2}\begin{bmatrix}0\\2\\-1\end{bmatrix}=\begin{bmatrix}(x+y)/2\\x-y\\y\end{bmatrix}.
\end{equation*}
Prescribing images of a dependent set requires consistency. There is no linear map with $T(1,0)=(1,1)$, $T(0,1)=(0,1)$ and $T(1,-1)=(2,2)$, since the first two conditions force $T(1,-1)=(1,0)$.
\end{Example}

\section{Matrices and Linear Transformations}\label{sec:Matrices}

\subsection{Column Space, Row Space and Null Space}

\begin{Definition}
For $A\in F^{m\times n}$, define $T_A\colon F^{n\times1}\to F^{m\times1}$ by $T_A(x)=Ax$. Its \newterm{column space} $\mathcal C(A)$ is the span of the columns of $A$, its \newterm{row space} $\mathcal R(A)$ is the span of its rows, and its \newterm{null space} is
\begin{equation*}
\mathcal N(A)=\qty{\, x\in F^{n\times1} \mid Ax=0 \,}.
\end{equation*}
The \newterm{column rank} and \newterm{row rank} are $\dim\mathcal C(A)$ and $\dim\mathcal R(A)$ respectively. The \newterm{nullity} of $A$ is $\dim\mathcal N(A)$.
\end{Definition}

Writing the columns as $a_1,\ldots,a_n$, we have
\begin{equation*}
Ax=\begin{bmatrix}a_1&\cdots&a_n\end{bmatrix}\begin{bmatrix}x_1\\\vdots\\x_n\end{bmatrix}=\sum_{j=1}^n x_ja_j.
\end{equation*}
Thus $\mathcal C(A)=\im T_A$ and $\mathcal N(A)=\ker T_A$. The row space is naturally identified, by transposition, with $\mathcal C(A^{\mathsf T})$. Also, if $B$ has columns $b_1,\ldots,b_p$, then $AB=\begin{bmatrix}Ab_1&\cdots&Ab_p\end{bmatrix}$.

\begin{Theorem}\label{thm:RowColumnRank}
The row rank and column rank of any matrix over any field are equal. Their common value is the \newterm{rank} of the matrix.
\end{Theorem}
\begin{proof}
If the column rank is $0$, the matrix is zero and so is its row rank. Otherwise, select $r$ columns forming a basis of $\mathcal C(A)$ and place them in a matrix $C\in F^{m\times r}$. Expressing every column of $A$ in this basis gives $A=CB$ for some $B\in F^{r\times n}$. Every row of $A$ is a linear combination of the $r$ rows of $B$. Hence the row rank of $A$ is at most its column rank. Applying the same argument to $A^{\mathsf T}$ gives the reverse inequality.
\end{proof}

\begin{Corollary}
For $A\in F^{m\times n}$,
\begin{equation*}
\rank A+\nullity A=n.
\end{equation*}
\end{Corollary}

\begin{Example}
For $A=\begin{bmatrix}1&2&0\\0&1&-1\end{bmatrix}$, the equations $Ax=0$ give $x_1=-2x_2$ and $x_3=x_2$. Therefore
\begin{equation*}
\mathcal N(A)=\operatorname{span}\qty{\begin{bmatrix}-2\\1\\1\end{bmatrix}},\qquad \rank A=2.
\end{equation*}
The matrices $\begin{bmatrix}1&2\\0&1\end{bmatrix}$ and $\begin{bmatrix}1&-1\\0&1\\1&1\end{bmatrix}$ both have rank $2$ and nullity $0$.
\end{Example}

\begin{Theorem}
If $A\in F^{m\times n}$ and $B\in F^{n\times p}$, then
\begin{equation*}
\rank(AB)\le\min\{\rank A,\rank B\},\qquad \nullity(AB)\ge\nullity B.
\end{equation*}
\end{Theorem}
\begin{proof}
We have $\mathcal C(AB)\subseteq\mathcal C(A)$ and $\mathcal N(B)\subseteq\mathcal N(AB)$, proving the first rank bound and the nullity bound. Moreover, $T_A$ maps $\mathcal C(B)$ onto $\mathcal C(AB)$, so rank--nullity applied to this restriction gives $\rank(AB)\le\rank B$.
\end{proof}

\subsection{Matrix Representation}

\begin{Definition}
Let $\mathcal B=(v_1,\ldots,v_n)$ and $\mathcal C=(w_1,\ldots,w_m)$ be ordered bases of $V$ and $W$. The \newterm{matrix representation} of $T\colon V\to W$ in these bases is
\begin{equation*}
[T]_{\mathcal C}^{\mathcal B}=\begin{bmatrix}[T(v_1)]_{\mathcal C}&\cdots&[T(v_n)]_{\mathcal C}\end{bmatrix}.
\end{equation*}
When $V=W$ and the same basis is used on both sides, write $[T]_{\mathcal B}$.
\end{Definition}

\begin{Theorem}\label{thm:MatrixCoordinates}
For every $v\in V$,
\begin{equation*}
[T(v)]_{\mathcal C}=[T]_{\mathcal C}^{\mathcal B}[v]_{\mathcal B}.
\end{equation*}
Conversely, every $m\times n$ matrix determines exactly one linear transformation in these bases.
\end{Theorem}
\begin{proof}
If $T(v_j)=\sum_i a_{ij}w_i$ and $v=\sum_j\alpha_jv_j$, then
\begin{equation*}
T(v)=\sum_i\pqty{\sum_j a_{ij}\alpha_j}w_i.
\end{equation*}
This is precisely the stated matrix equation. Conversely, the columns of any matrix prescribe the coordinates of $T(v_j)$; apply \cref{thm:SpecifyBasis}. In particular, every linear map between column-vector spaces is $T_A$ for a unique matrix $A$.
\end{proof}

\begin{Example}
For differentiation on $P_3(\mathbb R)$ with basis $(1,x,x^2,x^3)$,
\begin{equation*}
[D]=\begin{bmatrix}0&1&0&0\\0&0&2&0\\0&0&0&3\\0&0&0&0\end{bmatrix}.
\end{equation*}
For $D\colon P_4(\mathbb R)\to P_3(\mathbb R)$ and $J\colon P_3(\mathbb R)\to P_4(\mathbb R)$, where $Jp(x)=\int_0^x p(t)\,dt$, the standard matrices are
\begin{equation*}
A_D=\begin{bmatrix}0&1&0&0&0\\0&0&2&0&0\\0&0&0&3&0\\0&0&0&0&4\end{bmatrix},\qquad
A_J=\begin{bmatrix}0&0&0&0\\1&0&0&0\\0&1/2&0&0\\0&0&1/3&0\\0&0&0&1/4\end{bmatrix}.
\end{equation*}
Thus $A_DA_J=I_4$, whereas $A_JA_D=\operatorname{diag}(0,1,1,1,1)$. Indeed, $DJp=p$ but $JDp=p-p(0)$.
\end{Example}

\begin{Exercise}
On $P_3(\mathbb R)$, define $T(p)(x)=p(1-x)$. Find $[T]_{\mathcal B}$ for $\mathcal B=(1,x,x^2,x^3)$ by computing the coordinate columns of $T(1)$, $T(x)$, $T(x^2)$ and $T(x^3)$. Verify from the matrix that $T^2=\id$.
\end{Exercise}

\begin{Example}
Let $T(x,y)=(x-y,3y,2x+y)$, use the standard basis $\mathcal B$ of $\mathbb R^2$, and let
\begin{equation*}
\mathcal C=\pqty{\begin{bmatrix}1\\0\\1\end{bmatrix},\begin{bmatrix}1\\0\\-1\end{bmatrix},\begin{bmatrix}0\\1\\0\end{bmatrix}}.
\end{equation*}
Since
\begin{align*}
T(1,0)&=\tfrac32\begin{bmatrix}1\\0\\1\end{bmatrix}-\tfrac12\begin{bmatrix}1\\0\\-1\end{bmatrix},\\
T(0,1)&=-\begin{bmatrix}1\\0\\-1\end{bmatrix}+3\begin{bmatrix}0\\1\\0\end{bmatrix},
\end{align*}
we obtain
\begin{equation*}
[T]_{\mathcal C}^{\mathcal B}=\begin{bmatrix}3/2&0\\-1/2&-1\\0&3\end{bmatrix}.
\end{equation*}
\end{Example}

\subsection{Change of Basis}

\begin{Definition}
For ordered bases $\mathcal B$ and $\widetilde{\mathcal B}$ of $V$, the \newterm{change-of-basis matrix} from $\mathcal B$-coordinates to $\widetilde{\mathcal B}$-coordinates is
\begin{equation*}
P=[\id_V]_{\widetilde{\mathcal B}}^{\mathcal B},\qquad [v]_{\widetilde{\mathcal B}}=P[v]_{\mathcal B}.
\end{equation*}
It is invertible, with inverse $[\id_V]_{\mathcal B}^{\widetilde{\mathcal B}}$.
\end{Definition}

\begin{Theorem}\label{thm:ChangeBasis}
Suppose $A=[T]_{\mathcal C}^{\mathcal B}$ and $\widetilde A=[T]_{\widetilde{\mathcal C}}^{\widetilde{\mathcal B}}$. If $P$ changes domain coordinates from $\mathcal B$ to $\widetilde{\mathcal B}$ and $Q$ changes codomain coordinates from $\mathcal C$ to $\widetilde{\mathcal C}$, then
\begin{equation*}
\widetilde A=QAP^{-1}.
\end{equation*}
In particular, for an operator with the same basis change in domain and codomain, $[T]_{\widetilde{\mathcal B}}=P[T]_{\mathcal B}P^{-1}$.
\end{Theorem}
\begin{proof}
For all $v\in V$, $\widetilde A P[v]_{\mathcal B}=\widetilde A[v]_{\widetilde{\mathcal B}}=[T(v)]_{\widetilde{\mathcal C}}=QA[v]_{\mathcal B}$. Since coordinate vectors range over all of $F^n$, $\widetilde A P=QA$.
\end{proof}

\begin{Example}
For $T(x,y)=(y,x-y)$, the matrix in standard bases is $A=\begin{bmatrix}0&1\\1&-1\end{bmatrix}$. Use the domain and codomain bases
\begin{equation*}
\widetilde{\mathcal B}=\pqty{\begin{bmatrix}0\\1\end{bmatrix},\begin{bmatrix}1\\1\end{bmatrix}},\qquad \widetilde{\mathcal C}=\pqty{\begin{bmatrix}1\\-1\end{bmatrix},\begin{bmatrix}-1\\0\end{bmatrix}}.
\end{equation*}
Since $T(0,1)=(1,-1)$ and $T(1,1)=(1,0)$,
\begin{equation*}
\widetilde A=\begin{bmatrix}1&0\\0&-1\end{bmatrix}.
\end{equation*}
This uses different domain and codomain bases; it is not a diagonalisation of $T$.
\end{Example}

\section{Algebra of Linear Transformations}\label{sec:AlgebraTransformations}

\begin{Definition}
Write $L(V,W)$ for the set of all linear transformations from $V$ to $W$. Define addition and scalar multiplication pointwise:
\begin{equation*}
(S+T)(v)=S(v)+T(v),\qquad (\alpha T)(v)=\alpha T(v).
\end{equation*}
A linear transformation $V\to V$ is a \newterm{linear operator}; write $L(V)=L(V,V)$.
\end{Definition}

\begin{Theorem}
The set $L(V,W)$ is a vector space. If $\dim V=n$ and $\dim W=m$, then $\dim L(V,W)=mn$.
\end{Theorem}
\begin{proof}
Sums and scalar multiples of linear transformations remain linear, and the vector-space axioms hold pointwise. The zero transformation is the additive identity and $(-T)(v)=-T(v)$. In fixed bases, the correspondence $T\mapsto[T]_{\mathcal C}^{\mathcal B}$ is a linear bijection onto $F^{m\times n}$. More explicitly, the maps $E_{ij}$ specified by $E_{ij}(v_k)=\delta_{jk}w_i$ form a basis of $L(V,W)$, since any $T$ with matrix $[a_{ij}]$ has the unique expansion $T=\sum_{i,j}a_{ij}E_{ij}$.
\end{proof}

Composition of linear transformations is linear, associative and distributive over addition. For $T\colon U\to V$ and $S\colon V\to W$,
\begin{equation*}
[S\circ T]_{\mathcal D}^{\mathcal B}=[S]_{\mathcal D}^{\mathcal C}[T]_{\mathcal C}^{\mathcal B}.
\end{equation*}
Thus $L(V)$ is an associative algebra with identity $\id_V$, and in particular a ring under addition and composition. Set $T^0=\id_V$ and $T^{k+1}=T\circ T^k$.

\subsection{Isomorphisms}

\begin{Definition}
A linear transformation is \newterm{non-singular} if its kernel is $\{0\}$. It is an \newterm{isomorphism} if it has a linear inverse. In this case it is \newterm{invertible}, and the inverse is denoted by $T^{-1}$. Vector spaces are \newterm{isomorphic}, written $V\cong W$, if an isomorphism exists between them.
\end{Definition}

\begin{Theorem}
A linear transformation is an isomorphism if and only if it is bijective.
\end{Theorem}
\begin{proof}
An inverse implies bijectivity. Conversely, suppose $T$ is bijective. Write $w_1=T(v_1)$ and $w_2=T(v_2)$. Then
\begin{equation*}
T^{-1}(\alpha w_1+w_2)=T^{-1}T(\alpha v_1+v_2)=\alpha v_1+v_2=\alpha T^{-1}(w_1)+T^{-1}(w_2),
\end{equation*}
so the inverse function is linear.
\end{proof}

\begin{Theorem}\label{thm:IsomorphismCriteria}
If $\dim V=\dim W=n<\infty$, the following conditions on a linear map $T\colon V\to W$ are equivalent:
\begin{enumerate}[label=(\roman*)]
\item $T$ is an isomorphism;
\item $T$ is injective;
\item $T$ is surjective;
\item $T$ takes every basis of $V$ to a basis of $W$;
\item $T$ takes some basis of $V$ to a basis of $W$.
\end{enumerate}
\end{Theorem}
\begin{proof}
Rank--nullity gives $\nullity T=0$ if and only if $\rank T=n$, proving the equivalence of the first three conditions. An isomorphism preserves independence and spanning, so (i) implies (iv), and (iv) implies (v). If (v) holds, the image contains a basis of $W$, making $T$ surjective.
\end{proof}

\begin{Theorem}
Every $n$-dimensional vector space over $F$ is isomorphic to $F^n$. Consequently two finite dimensional vector spaces over the same field are isomorphic if and only if they have the same dimension.
\end{Theorem}
\begin{proof}
For an ordered basis $\mathcal B$, the coordinate map $v\mapsto[v]_{\mathcal B}$ is linear. Uniqueness of coordinates makes it injective, and every column of scalars is the coordinate vector of the corresponding linear combination, making it surjective. Conversely, an isomorphism takes a basis to a basis and therefore preserves dimension.
\end{proof}

\begin{Remark*}
Finite dimensionality is essential in \cref{thm:IsomorphismCriteria}. On $\mathbb R[x]$, differentiation is surjective but not injective, while multiplication by $x$ is injective but not surjective. Also, if $A\in F^{m\times n}$ and $B\in F^{n\times m}$ satisfy $AB=I_m$ and $BA=I_n$, the associated linear maps are inverses, so $m=n$.
\end{Remark*}


% !TeX root = Linear Algebra.tex
% Handwritten pages 11 and 21--23.

\section{Direct Sums and Projections}\label{sec:DirectSums}

\begin{Definition}
For subspaces $U,W\subseteq V$, their \newterm{sum} is
\begin{equation*}
U+W=\qty{\, u+w \mid u\in U,\ w\in W \,}.
\end{equation*}
The sum is \newterm{direct} if $U\cap W=\{0\}$; we then write $U\oplus W$. If $V=U\oplus W$, then $U$ is a \newterm{direct complement} of $W$ in $V$.
\end{Definition}

The sum is a subspace containing $U\cup W$, since $\alpha(u_1+w_1)+(u_2+w_2)=(\alpha u_1+u_2)+(\alpha w_1+w_2)$. Indeed, $U+W=\operatorname{span}(U\cup W)$.

\begin{Theorem}\label{thm:DirectUnique}
The equality $V=U\oplus W$ holds if and only if each $v\in V$ has exactly one expression $v=u+w$ with $u\in U$ and $w\in W$.
\end{Theorem}
\begin{proof}
If the sum is direct and $u+w=u'+w'$, then $u-u'=w'-w\in U\cap W$, so $u=u'$ and $w=w'$. Conversely, existence of the expressions gives $V=U+W$. For $z\in U\cap W$, the two expressions $z=z+0=0+z$ and uniqueness give $z=0$.
\end{proof}

\begin{Theorem}[Dimension Formula]\label{thm:DimensionSum}
For finite dimensional subspaces $U,W$,
\begin{equation*}
\dim(U+W)=\dim U+\dim W-\dim(U\cap W).
\end{equation*}
\end{Theorem}
\begin{proof}
Choose a basis $z_1,\ldots,z_k$ of $U\cap W$ and extend it to bases $(z_1,\ldots,z_k,u_1,\ldots,u_r)$ of $U$ and $(z_1,\ldots,z_k,w_1,\ldots,w_s)$ of $W$. The combined list $z_1,\ldots,z_k,u_1,\ldots,u_r,w_1,\ldots,w_s$ spans $U+W$. In any relation among these vectors, the sum of the $w_j$ terms lies in $U\cap W$, so independence of the basis of $W$ forces all their coefficients to vanish. Independence of the basis of $U$ then makes the remaining coefficients zero. Hence the combined list is a basis, and its length is $k+r+s=(k+r)+(k+s)-k$.
\end{proof}

\begin{Corollary}
Every subspace of a finite dimensional vector space has a direct complement.
\end{Corollary}
\begin{proof}
Extend a basis of $W$ to a basis of $V$ and let $U$ be the span of the added vectors. Uniqueness of basis coordinates gives $V=U\oplus W$.
\end{proof}

\begin{Remark*}
Complements need not be unique. In $\mathbb R^2$, the $x$-axis has both the $y$-axis and the line $y=x$ as direct complements.
\end{Remark*}

\subsection{Direct Products and Several Summands}

\begin{Definition}
For vector spaces $V,W$ over the same field, their \newterm{direct product} is
\begin{equation*}
V\times W=\qty{\, (v,w) \mid v\in V,\ w\in W \,},
\end{equation*}
with componentwise operations. For two factors this is also called the \newterm{external direct sum}.
\end{Definition}

The inclusions $i_V(v)=(v,0)$ and $i_W(w)=(0,w)$ and the projections $\pi_V(v,w)=v$, $\pi_W(v,w)=w$ are linear. They satisfy
\begin{equation*}
\pi_Vi_V=\id_V,\quad \pi_Wi_W=\id_W,\quad \pi_Vi_W=0,\quad \pi_Wi_V=0,\quad i_V\pi_V+i_W\pi_W=\id_{V\times W}.
\end{equation*}
In particular, $i_V\pi_V$ is a projection onto $V\times\{0\}$; it is not the identity on the entire product unless $W=\{0\}$. We have
\begin{equation*}
V\times W=(V\times\{0\})\oplus(\{0\}\times W),\qquad \dim(V\times W)=\dim V+\dim W
\end{equation*}
in finite dimensions, by taking the embedded bases of the factors. If $V=U\oplus W$ internally, the addition map $U\times W\to V$, $(u,w)\mapsto u+w$, is an isomorphism.

\begin{Exercise}
Given linear maps $T_V\colon Z\to V$ and $T_W\colon Z\to W$, show that there is exactly one linear map $T\colon Z\to V\times W$ satisfying $\pi_VT=T_V$ and $\pi_WT=T_W$. Explicitly, $T(z)=(T_V(z),T_W(z))$.
\end{Exercise}

\begin{Definition}
For finitely many subspaces $W_1,\ldots,W_k$, write $V=W_1\oplus\cdots\oplus W_k$ if every $v\in V$ has a unique decomposition $v=w_1+\cdots+w_k$ with $w_i\in W_i$.
\end{Definition}

Equivalently, the subspaces span $V$ and $W_i\cap\sum_{j\ne i}W_j=\{0\}$ for each $i$. Pairwise trivial intersections alone are not sufficient when there are three or more summands. Concatenating bases of the summands gives a basis of $V$, and conversely such a concatenated list is a basis precisely when the sum is direct and spans $V$. In particular, non-zero vectors $v_1,\ldots,v_n$ form a basis exactly when $V=\bigoplus_{i=1}^n\operatorname{span}\{v_i\}$.

\subsection{Idempotent Operators}

\begin{Definition}
An operator $P\in L(V)$ is \newterm{idempotent} if $P^2=P$. Such an operator is also called a \newterm{projection}.
\end{Definition}

\begin{Theorem}\label{thm:ProjectionSplit}
If $P^2=P$, then $V=\im P\oplus\ker P$. Conversely, a decomposition $V=U\oplus W$ defines a projection $P(u+w)=u$ with image $U$ and kernel $W$.
\end{Theorem}
\begin{proof}
For every $v$, write $v=P(v)+(v-P(v))$. The first term is in the image, and $P(v-P(v))=P(v)-P^2(v)=0$, so the second is in the kernel. If $z=P(u)\in\ker P$, then $z=P^2(u)=P(z)=0$. The sum is therefore direct. Conversely, uniqueness of the direct-sum decomposition makes $P$ well defined and linear, and applying it twice gives $P^2=P$. Its image and kernel are as claimed.
\end{proof}

\begin{Example}
On $\mathbb R^3$, the operators
\begin{equation*}
P(x,y,z)=(x,0,z),\qquad Q(x,y,z)=\pqty{\frac{x+y}{2},\frac{x+y}{2},0}
\end{equation*}
are idempotent. Here $\im P$ is the $xz$-plane and $\ker P$ the $y$-axis, while
\begin{equation*}
\im Q=\operatorname{span}\qty{\begin{bmatrix}1\\1\\0\end{bmatrix}},\qquad \ker Q=\qty{\, \begin{bmatrix}x\\-x\\z\end{bmatrix} \mid x,z\in\mathbb R \,}.
\end{equation*}
The zero and identity operators are also projections.
\end{Example}

\begin{Remark*}
For any operator $T$, $\im T\subseteq\ker T$ if and only if $T^2=0$. In contrast, $\im T\cap\ker T=\{0\}$ does not imply $T^2=T$. For example, $T=2\id_{\mathbb R}$ has trivial kernel but is not idempotent. A projection is distinguished by acting as the identity on its image.
\end{Remark*}

\begin{Exercise}
Assume that $\operatorname{char}F\ne2$. On $F^{n\times n}$ define $P(A)=\frac12(A+A^{\mathsf T})$. Show that $P$ is a projection, that its image is the space of symmetric matrices, and that its kernel is the space of skew-symmetric matrices. Deduce
\begin{equation*}
F^{n\times n}=\qty{\, A \mid A^{\mathsf T}=A \,}\oplus\qty{\, A \mid A^{\mathsf T}=-A \,}.
\end{equation*}
Explain why the hypothesis on the characteristic is necessary for this formula.
\end{Exercise}


% !TeX root = Linear Algebra.tex
% Handwritten pages 27--31. Page 28 is blank; prerequisites on annihilators are supplied.

\section{Linear Functionals and Dual Spaces}\label{sec:Duality}

\begin{Definition}
A \newterm{linear functional} on a vector space $V$ over $F$ is a linear map $V\to F$. The \newterm{algebraic dual} (or \newterm{dual space}) of $V$ is $V^*=L(V,F)$.
\end{Definition}

\begin{Example}
The map $f(x,y,z)=3x-y+2z$ is a linear functional on $\mathbb R^3$, whereas $g(x,y,z)=xyz$ is not. For $a\in F$, evaluation $\varepsilon_a\colon F[x]\to F$, $\varepsilon_a(p)=p(a)$, is linear. On $F^{n\times n}$, the trace $\operatorname{tr}(A)=\sum_i a_{ii}$ is linear. The determinant is not linear for $n\ge2$: for example, $\det E_{11}=\det(I_n-E_{11})=0$ but $\det I_n=1$. For $n=1$, the determinant is linear.
\end{Example}

\subsection{Dual Bases}

\begin{Theorem}\label{thm:DualBasis}
Let $\mathcal B=(v_1,\ldots,v_n)$ be an ordered basis of $V$. There are unique functionals $v_i^*$ satisfying $v_i^*(v_j)=\delta_{ij}$. They form a basis $\mathcal B^*=(v_1^*,\ldots,v_n^*)$ of $V^*$, called the \newterm{dual basis}, and
\begin{equation*}
f=\sum_{i=1}^n f(v_i)v_i^*\qquad(f\in V^*).
\end{equation*}
In particular, $\dim V^*=\dim V$.
\end{Theorem}
\begin{proof}
Define $v_i^*(\sum_j\alpha_jv_j)=\alpha_i$. This is linear and has the required values; uniqueness follows from \cref{thm:SpecifyBasis}. The expansion for $f$ holds by evaluating both sides on each $v_j$. If $\sum_i c_iv_i^*=0$, evaluation on $v_j$ gives $c_j=0$, proving independence.
\end{proof}

In coordinates, if $f(v_i)=c_i$, then
\begin{equation*}
f(v)=\begin{bmatrix}c_1&\cdots&c_n\end{bmatrix}[v]_{\mathcal B}.
\end{equation*}
The map $v_i\mapsto v_i^*$ extends to an isomorphism $V\to V^*$, but this isomorphism depends on the chosen basis.

\subsection{Dual Transformations}

\begin{Definition}
For a linear map $T\colon V\to W$, its \newterm{dual transformation} is
\begin{equation*}
T^*\colon W^*\to V^*,\qquad T^*(f)=f\circ T.
\end{equation*}
\end{Definition}

The composition $f\circ T$ is linear, and $T^*(\alpha f+g)=\alpha T^*(f)+T^*(g)$, so $T^*$ is linear. Notice that its direction is reversed.

\begin{Theorem}\label{thm:DualTranspose}
In finite dimensions, if $A=[T]_{\mathcal C}^{\mathcal B}$, then
\begin{equation*}
[T^*]_{\mathcal B^*}^{\mathcal C^*}=A^{\mathsf T}.
\end{equation*}
\end{Theorem}
\begin{proof}
Write $T(v_j)=\sum_i a_{ij}w_i$. Then $(T^*w_i^*)(v_j)=w_i^*(T(v_j))=a_{ij}$, so $T^*w_i^*=\sum_j a_{ij}v_j^*$. These are exactly the columns of $A^{\mathsf T}$.
\end{proof}

\begin{Example}
Let $T\colon\mathbb R^3\to\mathbb R^2$ be $T(x,y,z)=(y-2z,x+3y-z)$. If $f(s,t)=c_1s+c_2t$, then
\begin{equation*}
T^*f(x,y,z)=c_2x+(c_1+3c_2)y+(-2c_1-c_2)z.
\end{equation*}
Thus the standard matrices are
\begin{equation*}
[T]=\begin{bmatrix}0&1&-2\\1&3&-1\end{bmatrix},\qquad [T^*]=\begin{bmatrix}0&1\\1&3\\-2&-1\end{bmatrix}.
\end{equation*}
\end{Example}

\subsection{Annihilators and Hyperplanes}

\begin{Definition}
For $S\subseteq V$, its \newterm{annihilator} is the subspace
\begin{equation*}
S^\circ=\qty{\, f\in V^* \mid f(s)=0\text{ for every }s\in S \,}.
\end{equation*}
\end{Definition}

Linearity gives $S^\circ=(\operatorname{span}S)^\circ$. In particular, $\{0\}^\circ=V^*$ and $V^\circ=\{0\}$.

\begin{Theorem}\label{thm:AnnihilatorDimension}
If $\dim V=n$ and $W\subseteq V$ has dimension $k$, then $\dim W^\circ=n-k$.
\end{Theorem}
\begin{proof}
Extend a basis $v_1,\ldots,v_k$ of $W$ to a basis $v_1,\ldots,v_n$ of $V$. The expansion in \cref{thm:DualBasis} shows that a functional vanishes on $W$ precisely when its first $k$ dual-basis coefficients vanish. Therefore $(v_{k+1}^*,\ldots,v_n^*)$ is a basis of $W^\circ$.
\end{proof}

\begin{Definition}
A \newterm{hyperplane} (also called a \newterm{hyperspace}) of an $n$-dimensional vector space is a subspace of dimension $n-1$.
\end{Definition}

\begin{Theorem}
The hyperplanes of a finite dimensional vector space are exactly the kernels of its non-zero linear functionals. Every $k$-dimensional subspace of an $n$-dimensional space is an intersection of $n-k$ hyperplanes.
\end{Theorem}
\begin{proof}
A non-zero functional has image $F$ and hence rank $1$, so its kernel has dimension $n-1$. Conversely, for a hyperplane $W$, extend its basis by one vector and take the last dual-basis functional: its kernel is $W$. More generally, with the basis in \cref{thm:AnnihilatorDimension},
\begin{equation*}
W=\bigcap_{i=k+1}^n\ker v_i^*=\bigcap_{f\in W^\circ}\ker f.
\end{equation*}
When $W=V$, the empty intersection is understood as $V$.
\end{proof}

\begin{Theorem}
If $T\colon V\to W$ is linear, then $\ker T^*=(\im T)^\circ$. If $W$ is finite dimensional, then $\rank T^*=\rank T$.
\end{Theorem}
\begin{proof}
For $f\in W^*$, $T^*f=0$ if and only if $f(T(v))=0$ for every $v\in V$, which is equivalent to $f\in(\im T)^\circ$. If $\dim W=m<\infty$, rank--nullity and the annihilator dimension formula give
\begin{equation*}
\rank T^*=m-\dim(\im T)^\circ=m-(m-\dim\im T)=\rank T.
\end{equation*}
\end{proof}

\begin{Remark*}
Together with \cref{thm:DualTranspose}, this gives a second proof of equality of row and column rank. Unlike orthogonality arguments involving $A^{\mathsf T}A$, both proofs work over an arbitrary field.
\end{Remark*}

\begin{Theorem}
For subspaces $U,W$ of a finite dimensional space $V$,
\begin{align*}
U\subseteq W&\iff W^\circ\subseteq U^\circ,\\
(U+W)^\circ&=U^\circ\cap W^\circ,\\
(U\cap W)^\circ&=U^\circ+W^\circ.
\end{align*}
\end{Theorem}
\begin{proof}
The forward inclusion in the first assertion follows from the definition, and the reverse follows from the description of a subspace as the intersection of the kernels of its annihilator. A functional vanishes on $U+W$ exactly when it vanishes on both $U$ and $W$, proving the second assertion. For the third, $U^\circ+W^\circ\subseteq(U\cap W)^\circ$ is immediate. Using the second assertion, \cref{thm:DimensionSum,thm:AnnihilatorDimension} give
\begin{align*}
\dim(U^\circ+W^\circ)&=(n-\dim U)+(n-\dim W)-\bigl(n-\dim(U+W)\bigr)\\
&=n-\dim(U\cap W)=\dim(U\cap W)^\circ.
\end{align*}
The inclusion is therefore equality.
\end{proof}

\subsection{The Double Dual}

\begin{Definition}
The \newterm{double dual} of $V$ is $V^{**}=(V^*)^*$. For each $v\in V$, define $\varepsilon_v\in V^{**}$ by $\varepsilon_v(f)=f(v)$.
\end{Definition}

\begin{Theorem}
If $V$ is finite dimensional, the evaluation map $J\colon V\to V^{**}$, $J(v)=\varepsilon_v$, is an isomorphism which does not depend on a choice of basis.
\end{Theorem}
\begin{proof}
For $f,g\in V^*$, $\varepsilon_v(\alpha f+g)=\alpha f(v)+g(v)$, so $\varepsilon_v$ is a functional. For $u,v\in V$, evaluation on any $f$ gives $J(\alpha u+v)(f)=\alpha f(u)+f(v)$, proving linearity of $J$. If $v\ne0$, extend $v$ to a basis and take the first dual-basis functional $f$; then $J(v)(f)=1$. Thus $J$ is injective. Since $\dim V^{**}=\dim V$, it is also surjective.
\end{proof}


% !TeX root = Linear Algebra.tex
% Handwritten pages 31--41.

\section{Invariant Subspaces and Eigenvalues}\label{sec:Eigenvalues}

\begin{Definition}
A subspace $W\subseteq V$ is \newterm{invariant} under $T\in L(V)$ (or \newterm{$T$-invariant}) if $T(W)\subseteq W$. The restriction $T|_W$ is then an operator on $W$.
\end{Definition}

\begin{Example}
The subspaces $\{0\}$, $V$, $\ker T$ and $\im T$ are always invariant under $T$. So is every subspace contained in $\ker T$ or containing $\im T$. For a scalar operator $T=\alpha\id_V$, every subspace is invariant.

For $T(x,y)=(x+y,x-y)$ on $\mathbb R^2$, the $x$-axis is not invariant, since $T(1,0)=(1,1)$. However, the lines spanned by $\begin{bmatrix}1\\\sqrt2-1\end{bmatrix}$ and $\begin{bmatrix}1\\-\sqrt2-1\end{bmatrix}$ are invariant: $T$ acts on them by multiplication by $\sqrt2$ and $-\sqrt2$ respectively.
\end{Example}

\begin{Theorem}\label{thm:InvariantBlocks}
Let $W$ be a subspace of a finite dimensional $V$, and extend an ordered basis of $W$ to an ordered basis $\mathcal B$ of $V$. Then $W$ is invariant under $T$ if and only if
\begin{equation*}
[T]_{\mathcal B}=\begin{bmatrix}A&B\\0&C\end{bmatrix},
\end{equation*}
where $A$ represents $T|_W$. If $V=W_1\oplus W_2$ and the basis is the concatenation of bases of $W_1,W_2$, both summands are invariant if and only if the matrix is block diagonal:
\begin{equation*}
[T]_{\mathcal B}=\begin{bmatrix}[T|_{W_1}]_{\mathcal B_1}&0\\0&[T|_{W_2}]_{\mathcal B_2}\end{bmatrix}.
\end{equation*}
\end{Theorem}
\begin{proof}
The columns corresponding to basis vectors of $W$ have zero coordinates outside $W$ exactly when their images lie in $W$. This is the first block condition. Applying the same observation to the columns for each summand gives the second.
\end{proof}

\subsection{Eigenvalues and Eigenvectors}

\begin{Definition}
A scalar $\lambda\in F$ is an \newterm{eigenvalue} of $T\in L(V)$ if there is a non-zero vector $v\in V$ with $T(v)=\lambda v$. Such a vector is an \newterm{eigenvector}, and $(\lambda,v)$ is an \newterm{eigenpair}. The \newterm{eigenspace} corresponding to $\lambda$ is
\begin{equation*}
E(\lambda,T)=\ker(T-\lambda\id_V).
\end{equation*}
The same terminology applies to a square matrix $A$ when $Av=\lambda v$ with $v\ne0$.
\end{Definition}

A non-zero vector $v$ is an eigenvector exactly when $\operatorname{span}\{v\}$ is invariant. Every non-zero scalar multiple of an eigenvector is an eigenvector for the same eigenvalue. The eigenspace contains the zero vector, although the zero vector is never an eigenvector.

\begin{Example}
The real operator $T(x,y)=(-y,x)$ has no eigenvalues over $\mathbb R$: the equations $-y=\lambda x$ and $x=\lambda y$ force $\lambda^2+1=0$ for a non-zero solution. Over $\mathbb C$, its eigenpairs include
\begin{equation*}
\pqty{i,\begin{bmatrix}1\\-i\end{bmatrix}},\qquad \pqty{-i,\begin{bmatrix}1\\i\end{bmatrix}}.
\end{equation*}
Thus eigenvalues depend on the scalar field.
\end{Example}

\begin{Example}
For $T(x,y)=(x+y,x)$ on $\mathbb R^2$, the eigenvalue equation yields $\lambda^2-\lambda-1=0$. The eigenvalues are $\lambda_\pm=(1\pm\sqrt5)/2$, with eigenvectors $\begin{bmatrix}1\\\lambda_\pm-1\end{bmatrix}$. For $S(x,y)=(x+y,y)$, the only eigenvalue is $1$, and $E(1,S)=\operatorname{span}\qty{\begin{bmatrix}1\\0\end{bmatrix}}$. Hence $S$ has no basis of eigenvectors.
\end{Example}

\begin{Theorem}\label{thm:DistinctEigenvalues}
Eigenvectors corresponding to distinct eigenvalues are linearly independent. More generally, the sum of the eigenspaces for distinct eigenvalues is direct.
\end{Theorem}
\begin{proof}
Suppose $v_1,\ldots,v_k$ is a dependent list of eigenvectors for distinct eigenvalues, with $k$ as small as possible. In a non-trivial relation $\sum_{i=1}^k\alpha_iv_i=0$, all coefficients are non-zero. Apply $T-\lambda_k\id_V$ to obtain
\begin{equation*}
\sum_{i=1}^{k-1}\alpha_i(\lambda_i-\lambda_k)v_i=0.
\end{equation*}
All displayed coefficients are non-zero, contradicting minimality. For the statement about eigenspaces, a non-zero relation among vectors from distinct eigenspaces would give exactly such a dependence after omitting zero terms.
\end{proof}

\begin{Corollary}
An operator on an $n$-dimensional space has at most $n$ distinct eigenvalues.
\end{Corollary}

\subsection{Characteristic Polynomials}

\begin{Definition}
Square matrices $A,B\in F^{n\times n}$ are \newterm{similar} if $B=PAP^{-1}$ for some invertible $P$. The \newterm{characteristic polynomial} of $A$ is $\chi_A(x)=\det(xI_n-A)$. The equation $\chi_A(x)=0$ is its \newterm{characteristic equation}, and its roots are its \newterm{characteristic roots}.
\end{Definition}

\begin{Theorem}
Similar matrices have the same characteristic polynomial. Thus $\chi_T=\chi_{[T]_{\mathcal B}}$ is independent of the basis $\mathcal B$. For $\lambda\in F$, the following are equivalent:
\begin{enumerate}[label=(\roman*)]
\item $\lambda$ is an eigenvalue of $T$;
\item $T-\lambda\id_V$ is not injective;
\item $T-\lambda\id_V$ is not surjective;
\item $T-\lambda\id_V$ is not invertible;
\item $\chi_T(\lambda)=0$.
\end{enumerate}
\end{Theorem}
\begin{proof}
For $B=PAP^{-1}$, $xI-B=P(xI-A)P^{-1}$, whose determinant is $\det(xI-A)$. The first two conditions express the existence of a non-zero vector in the kernel. The next two are equivalent by \cref{thm:IsomorphismCriteria}; the last uses the determinant criterion for invertibility of a square matrix.
\end{proof}

By the factor theorem, $p(\lambda)=0$ if and only if $(x-\lambda)$ divides $p(x)$. A non-zero polynomial of degree $n$ over any field has at most $n$ distinct roots. The fundamental theorem of algebra states that every non-constant complex polynomial has a complex root, and consequently splits completely into linear factors over $\mathbb C$.

\begin{Example}\label{ex:ThreeByThreeEigenvalues}
For
\begin{equation*}
A=\begin{bmatrix}1&-1&-1\\-1&1&-1\\-1&-1&1\end{bmatrix},
\end{equation*}
expanding $\det(xI-A)$ gives
\begin{equation*}
\chi_A(x)=x^3-3x^2+4=(x+1)(x-2)^2.
\end{equation*}
Over $\mathbb R$,
\begin{equation*}
E(-1,A)=\operatorname{span}\qty{\begin{bmatrix}1\\1\\1\end{bmatrix}},\qquad E(2,A)=\qty{\, \begin{bmatrix}x\\y\\z\end{bmatrix} \mid x+y+z=0 \,}.
\end{equation*}
\end{Example}

\begin{Remark*}
For an $n\times n$ matrix,
\begin{equation*}
\chi_A(x)=x^n-c_1x^{n-1}+c_2x^{n-2}-\cdots+(-1)^nc_n,
\end{equation*}
where $c_k$ is the sum of the principal minors of order $k$. A principal minor is obtained by choosing the same set of row and column indices. In particular, $c_1=\operatorname{tr}A$ and $c_n=\det A$. This formula follows from the determinant expansion: to contribute a term of degree $n-k$, choose the $x$ term at $n-k$ diagonal positions and the entries of $-A$ on the remaining principal submatrix.
\end{Remark*}

\section{Polynomials of Linear Operators}\label{sec:OperatorPolynomials}

\subsection{Polynomial Evaluation and the Minimal Polynomial}

For $p(x)=a_0+a_1x+\cdots+a_dx^d$ and $T\in L(V)$, define
\begin{equation*}
p(T)=a_0\id_V+a_1T+\cdots+a_dT^d.
\end{equation*}
The evaluation map $F[x]\to L(V)$, $p\mapsto p(T)$, preserves addition, multiplication, scalar multiplication and identity. In particular, polynomials in $T$ commute with each other and with $T$.

\begin{Theorem}\label{thm:PolynomialEigenvector}
If $T(v)=\lambda v$ with $v\ne0$, then $p(T)v=p(\lambda)v$. Thus $(p(\lambda),v)$ is an eigenpair of $p(T)$. If $T$ is invertible, then $\lambda\ne0$ and $T^{-1}v=\lambda^{-1}v$.
\end{Theorem}
\begin{proof}
Induction gives $T^jv=\lambda^jv$, and the polynomial identity follows by taking the indicated linear combination. If $T$ is invertible, an eigenvalue cannot be zero; apply $T^{-1}$ to $T(v)=\lambda v$.
\end{proof}

\begin{Definition}
A polynomial $p\in F[x]$ \newterm{annihilates} $T$ if $p(T)=0$. For a finite dimensional space, the \newterm{minimal polynomial} $m_T$ is the unique monic polynomial of least degree which annihilates $T$.
\end{Definition}

\begin{Theorem}\label{thm:MinimalPolynomial}
The minimal polynomial exists and is unique. A polynomial annihilates $T$ if and only if it is divisible by $m_T$.
\end{Theorem}
\begin{proof}
If $n=\dim V>0$, the $n^2+1$ operators $\id_V,T,\ldots,T^{n^2}$ are linearly dependent in the $n^2$-dimensional space $L(V)$. A non-zero annihilating polynomial therefore exists; divide by its leading coefficient and choose one of least degree. If two monic polynomials of this degree annihilate $T$, their difference has smaller degree and also annihilates $T$, so the difference must be zero.

For any annihilating polynomial $p$, polynomial division gives $p=qm_T+r$, where $r=0$ or $\deg r<\deg m_T$. Evaluating at $T$ gives $r(T)=0$, so minimality forces $r=0$. Conversely, every multiple of $m_T$ annihilates $T$. If $V=\{0\}$, take $m_T=1$; the assertions still hold.
\end{proof}

Equivalently, the kernel of the evaluation map is the principal ideal $(m_T)$ in $F[x]$.

\begin{Theorem}
The eigenvalues of $T$ in $F$ are exactly the roots of $m_T$ in $F$.
\end{Theorem}
\begin{proof}
If $T(v)=\lambda v$ with $v\ne0$, then $0=m_T(T)v=m_T(\lambda)v$, so $m_T(\lambda)=0$. Conversely, if $m_T(x)=(x-\lambda)q(x)$, minimality ensures $q(T)\ne0$. Choose $u$ with $v=q(T)u\ne0$. Then $(T-\lambda\id_V)v=m_T(T)u=0$, so $\lambda$ is an eigenvalue.
\end{proof}

\subsection{Cyclic Subspaces}

\begin{Definition}
For $v\in V$, the \newterm{$T$-cyclic subspace} generated by $v$ is
\begin{equation*}
Z(v,T)=\operatorname{span}\{v,Tv,T^2v,\ldots\}=\qty{\, p(T)v \mid p\in F[x] \,}.
\end{equation*}
\end{Definition}

This is the smallest $T$-invariant subspace containing $v$: it is invariant by construction, and every invariant subspace containing $v$ contains all its iterates.

\begin{Theorem}\label{thm:CyclicBasis}
In finite dimensions, there is a unique monic polynomial $m_v$ of least degree such that $m_v(T)v=0$. If $v\ne0$ and $\deg m_v=k$, then $(v,Tv,\ldots,T^{k-1}v)$ is a basis of $Z(v,T)$. In particular, $\dim Z(v,T)=\deg m_v$. For $v=0$, take $m_v=1$.
\end{Theorem}
\begin{proof}
The iterates eventually become dependent. Let $k$ be the first index for which the list
\begin{equation*}
v,Tv,\ldots,T^kv
\end{equation*}
is dependent. The first $k$ vectors are independent and
\begin{equation*}
T^kv=-a_0v-a_1Tv-\cdots-a_{k-1}T^{k-1}v.
\end{equation*}
Their span is invariant under $T$, so it contains every later iterate and is exactly $Z(v,T)$. The polynomial $x^k+a_{k-1}x^{k-1}+\cdots+a_0$ has the required minimal degree, and uniqueness follows by subtracting two monic candidates. The same polynomial-division argument as before shows that $m_v$ divides every $p$ with $p(T)v=0$.
\end{proof}

\begin{Example}
Let $A=\begin{bmatrix}1&0&-1\\0&2&1\\0&0&2\end{bmatrix}$ on $\mathbb R^3$, and let $v=\begin{bmatrix}1\\1\\1\end{bmatrix}$. Then
\begin{equation*}
Av=\begin{bmatrix}0\\3\\2\end{bmatrix},\qquad A^2v=\begin{bmatrix}-2\\8\\4\end{bmatrix},\qquad \det\begin{bmatrix}v&Av&A^2v\end{bmatrix}=-6.
\end{equation*}
Thus $Z(v,A)=\mathbb R^3$. In contrast, for $w=\begin{bmatrix}0\\1\\0\end{bmatrix}$, $Aw=2w$, so $Z(w,A)=\operatorname{span}\{w\}$ and $m_w(x)=x-2$.
\end{Example}

\begin{Theorem}\label{thm:MinimalDegreeBound}
If $\dim V=n$, then $\deg m_T\le n$.
\end{Theorem}
\begin{proof}
Induct on $n$, the case $n=0$ being immediate. Choose $v\ne0$ and set $k=\dim Z(v,T)=\deg m_v$. Let $U=\im m_v(T)$. Since $T$ commutes with $m_v(T)$, $U$ is invariant. Moreover, $Z(v,T)\subseteq\ker m_v(T)$, so rank--nullity gives $\dim U\le n-k<n$. By induction the minimal polynomial $q$ of $T|_U$ has degree at most $n-k$. For every $w\in V$, $m_v(T)w\in U$, so $q(T)m_v(T)w=0$. Consequently $qm_v$ annihilates $T$ and
\begin{equation*}
\deg m_T\le\deg q+\deg m_v\le(n-k)+k=n.
\end{equation*}
\end{proof}

\subsection{Companion Matrices and Cayley--Hamilton}

\begin{Definition}
For a monic polynomial $p(x)=x^k+a_{k-1}x^{k-1}+\cdots+a_1x+a_0$ of positive degree, its \newterm{companion matrix} is
\begin{equation*}
C(p)=\begin{bmatrix}
0&0&\cdots&0&-a_0\\
1&0&\cdots&0&-a_1\\
0&1&\cdots&0&-a_2\\
\vdots&\vdots&\ddots&\vdots&\vdots\\
0&0&\cdots&1&-a_{k-1}
\end{bmatrix}.
\end{equation*}
For $k=1$, this means $C(p)=\begin{bmatrix}-a_0\end{bmatrix}$.
\end{Definition}

\begin{Theorem}\label{thm:Companion}
Both the minimal and characteristic polynomials of $C(p)$ are $p$.
\end{Theorem}
\begin{proof}
Write $C=C(p)$. Since $e_1,Ce_1,\ldots,C^{k-1}e_1$ are $e_1,\ldots,e_k$, no polynomial of degree less than $k$ annihilates $e_1$. The last column gives $p(C)e_1=0$. As $p(C)$ commutes with $C$, it also kills $C^je_1=e_{j+1}$ for $0\le j<k$, and hence $p(C)=0$. Thus $m_C=p$.

For the characteristic polynomial, expand $\det(xI-C)$ along the last column. For a row $i<k$, the entry $a_{i-1}$ has cofactor $x^{i-1}$: the minor has determinant $x^{i-1}(-1)^{k-i}$ and the cofactor sign cancels this sign. The last entry $x+a_{k-1}$ has cofactor $x^{k-1}$. Therefore
\begin{equation*}
\det(xI-C)=a_0+a_1x+\cdots+a_{k-2}x^{k-2}+(x+a_{k-1})x^{k-1}=p(x).
\end{equation*}
The case $k=1$ is immediate.
\end{proof}

\begin{Theorem}[Cayley--Hamilton Theorem]\label{thm:CayleyHamilton}
Every operator on a finite dimensional vector space is annihilated by its characteristic polynomial:
\begin{equation*}
\chi_T(T)=0.
\end{equation*}
Consequently $m_T$ divides $\chi_T$.
\end{Theorem}
\begin{proof}
Fix $v\ne0$ and let $W=Z(v,T)$. In the cyclic basis from \cref{thm:CyclicBasis}, the matrix of $T|_W$ is $C(m_v)$. Extend this basis to one of $V$. By invariance,
\begin{equation*}
[T]=\begin{bmatrix}C(m_v)&B\\0&M\end{bmatrix},\qquad \chi_T(x)=\det(xI-C(m_v))\det(xI-M)=m_v(x)q(x).
\end{equation*}
Here we used \cref{thm:Companion} and the determinant of a block upper triangular matrix. Hence $\chi_T(T)v=q(T)m_v(T)v=0$. This holds for every non-zero $v$, and trivially for $v=0$, proving the operator identity. Divisibility follows from \cref{thm:MinimalPolynomial}.
\end{proof}

\section{Diagonalisation}\label{sec:Diagonalisation}

\begin{Definition}
An operator $T$ on a finite dimensional space is \newterm{diagonalisable} over $F$ if some ordered basis gives a diagonal matrix $[T]_{\mathcal B}$. A square matrix is diagonalisable if it is similar over $F$ to a diagonal matrix.
\end{Definition}

The columns of a diagonal matrix show that $T$ is diagonalisable exactly when $V$ has a basis of eigenvectors. Equivalently, $V$ is a direct sum of one-dimensional invariant subspaces.

\begin{Definition}
The \newterm{geometric multiplicity} of an eigenvalue $\lambda$ is $\dim E(\lambda,T)$. Its \newterm{algebraic multiplicity} is its multiplicity as a root of $\chi_T$, namely the largest $a$ such that $(x-\lambda)^a$ divides $\chi_T(x)$.
\end{Definition}

\begin{Theorem}\label{thm:MultiplicityBound}
For every eigenvalue, geometric multiplicity is at least $1$ and at most algebraic multiplicity.
\end{Theorem}
\begin{proof}
Let $k=\dim E(\lambda,T)\ge1$. Extend a basis of the eigenspace to one of $V$. Its invariance gives
\begin{equation*}
[T]=\begin{bmatrix}\lambda I_k&B\\0&M\end{bmatrix},\qquad \chi_T(x)=(x-\lambda)^k\det(xI-M).
\end{equation*}
Thus the algebraic multiplicity is at least $k$.
\end{proof}

\begin{Theorem}\label{thm:DiagonalisationCriteria}
For an operator $T$ on a finite dimensional vector space over $F$, the following are equivalent:
\begin{enumerate}[label=(\roman*)]
\item $T$ is diagonalisable over $F$;
\item $V$ is the direct sum of its eigenspaces;
\item $\chi_T$ splits into linear factors over $F$, and the geometric and algebraic multiplicities agree for each eigenvalue;
\item $m_T$ splits into distinct linear factors over $F$.
\end{enumerate}
\end{Theorem}
\begin{proof}
A basis of eigenvectors, grouped according to eigenvalue, gives (ii), and conversely concatenating bases of the eigenspaces gives such a basis. In a diagonal matrix, the number of occurrences of $\lambda$ on the diagonal is both its geometric and algebraic multiplicity, giving (i)$\Rightarrow$(iii). If (iii) holds, the sum of the eigenspace dimensions is $\deg\chi_T=\dim V$, and the sum is direct by \cref{thm:DistinctEigenvalues}; hence (ii) holds.

If $T$ is diagonal in an eigenbasis and its distinct eigenvalues are $\lambda_1,\ldots,\lambda_s$, then $\prod_i(T-\lambda_i\id_V)=0$. Every annihilating polynomial must vanish at every $\lambda_i$, so $m_T(x)=\prod_i(x-\lambda_i)$, proving (iv).

Conversely, suppose $m_T(x)=\prod_{i=1}^s(x-\lambda_i)$ with distinct $\lambda_i\in F$. Define
\begin{equation*}
q_i(x)=\prod_{j\ne i}\frac{x-\lambda_j}{\lambda_i-\lambda_j}.
\end{equation*}
The polynomial $\sum_iq_i-1$ has degree less than $s$ and vanishes at every $\lambda_i$, so it is zero. Also, $(x-\lambda_i)q_i(x)$ is a scalar multiple of $m_T(x)$, so $q_i(T)v\in E(\lambda_i,T)$. Thus $v=\sum_iq_i(T)v$ lies in the sum of the eigenspaces, which is direct by \cref{thm:DistinctEigenvalues}. This proves (ii). For the zero space, use $m_T=\chi_T=1$ and the empty direct sum.
\end{proof}

\begin{Corollary}
An operator on an $n$-dimensional space with $n$ distinct eigenvalues in $F$ is diagonalisable over $F$.
\end{Corollary}

\begin{Example}
The matrix in \cref{ex:ThreeByThreeEigenvalues} is diagonalisable over $\mathbb R$: the eigenbasis
\begin{equation*}
\pqty{\begin{bmatrix}1\\1\\1\end{bmatrix},\begin{bmatrix}1\\-1\\0\end{bmatrix},\begin{bmatrix}1\\0\\-1\end{bmatrix}}
\end{equation*}
gives $\operatorname{diag}(-1,2,2)$. Its minimal polynomial is $(x+1)(x-2)$. In contrast,
\begin{equation*}
\begin{bmatrix}1&1\\0&1\end{bmatrix}
\end{equation*}
has characteristic and minimal polynomial $(x-1)^2$ and is not diagonalisable. The real rotation
\begin{equation*}
\begin{bmatrix}0&-1\\1&0\end{bmatrix}
\end{equation*}
is diagonalisable over $\mathbb C$ but not over $\mathbb R$.
\end{Example}

\begin{Remark*}
Over $\mathbb C$, the splitting condition in (iii) is automatic by the fundamental theorem of algebra. Over a general field it must be stated explicitly. Having equal multiplicities for all eigenvalues that happen to lie in $F$ is not enough if the characteristic polynomial does not split over $F$.
\end{Remark*}


% !TeX root = Linear Algebra.tex
% Handwritten pages 41--46.

\section{Inner Product Spaces}\label{sec:InnerProducts}

Throughout this section, $F$ is either $\mathbb R$ or $\mathbb C$.

\begin{Definition}
An \newterm{inner product} on a vector space $V$ over $F$ is a function $\langle\cdot,\cdot\rangle\colon V\times V\to F$ such that, for all $u,v,w\in V$ and $\alpha\in F$,
\begin{align*}
\langle u+v,w\rangle&=\langle u,w\rangle+\langle v,w\rangle,\\
\langle\alpha u,v\rangle&=\alpha\langle u,v\rangle,\\
\langle u,v\rangle&=\overline{\langle v,u\rangle},\\
\langle v,v\rangle&\ge0,
\end{align*}
and $\langle v,v\rangle=0$ if and only if $v=0$. A vector space equipped with an inner product is an \newterm{inner product space}.
\end{Definition}

Thus the first argument is linear, while the second is conjugate linear:
\begin{equation*}
\langle u,\alpha v+w\rangle=\overline\alpha\langle u,v\rangle+\langle u,w\rangle.
\end{equation*}
Over $\mathbb R$, conjugation has no effect and the inner product is bilinear and symmetric.

\begin{Example}
The \newterm{standard inner product} on $F^{n\times1}$ is
\begin{equation*}
\langle u,v\rangle=v^*u=\sum_{i=1}^n u_i\overline{v_i},
\end{equation*}
where $v^*=\overline v^{\mathsf T}$. Over $\mathbb R$, this is the usual dot product $v^{\mathsf T}u$. On $F^{m\times n}$, the Frobenius inner product is
\begin{equation*}
\langle A,B\rangle=\operatorname{tr}(B^*A)=\sum_{i,j}a_{ij}\overline{b_{ij}}.
\end{equation*}
\end{Example}

\begin{Definition}
The \newterm{norm induced by an inner product} is
\begin{equation*}
\norm{v}=\sqrt{\langle v,v\rangle}.
\end{equation*}
Two vectors are \newterm{orthogonal}, written $u\perp v$, if $\langle u,v\rangle=0$.
\end{Definition}

The induced norm satisfies $\norm{v}\ge0$, with equality precisely for $v=0$, and
\begin{equation*}
\norm{\alpha v}=\abs{\alpha}\norm{v}.
\end{equation*}
The triangle inequality will follow from Cauchy--Schwarz.

\begin{Theorem}[Pythagoras]\label{thm:Pythagoras}
If $u\perp v$, then
\begin{equation*}
\norm{u+v}^2=\norm{u}^2+\norm{v}^2.
\end{equation*}
\end{Theorem}
\begin{proof}
Expand the left-hand side:
\begin{equation*}
\langle u+v,u+v\rangle=\norm{u}^2+\langle u,v\rangle+\overline{\langle u,v\rangle}+\norm{v}^2.
\end{equation*}
The middle terms vanish.
\end{proof}

\begin{Theorem}[Cauchy--Schwarz Inequality]\label{thm:CauchySchwarz}
For all $u,v\in V$,
\begin{equation*}
\abs{\langle u,v\rangle}\le\norm{u}\norm{v}.
\end{equation*}
Equality holds if and only if $u$ and $v$ are linearly dependent.
\end{Theorem}
\begin{proof}
If $v=0$, the assertion is immediate. Otherwise set
\begin{equation*}
\alpha=\frac{\langle u,v\rangle}{\norm{v}^2}.
\end{equation*}
Then $u-\alpha v\perp v$, and Pythagoras applied to $u=(u-\alpha v)+\alpha v$ gives
\begin{equation*}
\norm{u}^2=\norm{u-\alpha v}^2+\frac{\abs{\langle u,v\rangle}^2}{\norm{v}^2}\ge\frac{\abs{\langle u,v\rangle}^2}{\norm{v}^2}.
\end{equation*}
Taking square roots proves the inequality. Equality holds precisely when $u-\alpha v=0$, which means that $u$ and $v$ are dependent. The case $v=0$ also satisfies this characterisation.
\end{proof}

\begin{Corollary}[Triangle Inequality]
For all $u,v\in V$,
\begin{equation*}
\norm{u+v}\le\norm{u}+\norm{v}.
\end{equation*}
\end{Corollary}
\begin{proof}
By Cauchy--Schwarz,
\begin{align*}
\norm{u+v}^2&=\norm{u}^2+2\Re\langle u,v\rangle+\norm{v}^2\\
&\le\norm{u}^2+2\norm{u}\norm{v}+\norm{v}^2=(\norm{u}+\norm{v})^2.
\end{align*}
\end{proof}

\subsection{Orthogonal and Orthonormal Sets}

\begin{Definition}
A set $S\subseteq V$ is \newterm{orthogonal} if distinct vectors in $S$ are orthogonal. It is \newterm{orthonormal} if it is orthogonal and every vector in it has norm $1$. An \newterm{orthonormal basis} is a basis which is orthonormal.
\end{Definition}

Every orthogonal set of non-zero vectors is linearly independent. Indeed, if $\sum_i\alpha_iv_i=0$, taking the inner product with $v_j$ gives $\alpha_j\norm{v_j}^2=0$. In particular, every orthonormal set is independent.

\begin{Example}
In $\mathbb R^3$, the vectors $\begin{bmatrix}1\\0\\1\end{bmatrix}$ and $\begin{bmatrix}1\\0\\-1\end{bmatrix}$ are orthogonal but not orthonormal. Dividing both by $\sqrt2$ gives an orthonormal set. In $\mathbb C^2$, the vectors $\begin{bmatrix}1\\i\end{bmatrix}$ and $\begin{bmatrix}1\\-i\end{bmatrix}$ are orthogonal because
\begin{equation*}
\left\langle\begin{bmatrix}1\\i\end{bmatrix},\begin{bmatrix}1\\-i\end{bmatrix}\right\rangle=1+i\overline{-i}=0.
\end{equation*}
\end{Example}

\begin{Theorem}\label{thm:ONCoordinates}
If $\mathcal B=(u_1,\ldots,u_n)$ is an orthonormal basis of $V$, then
\begin{equation*}
v=\sum_{i=1}^n\langle v,u_i\rangle u_i,
\qquad
[v]_{\mathcal B}=\begin{bmatrix}\langle v,u_1\rangle\\\vdots\\\langle v,u_n\rangle\end{bmatrix},
\qquad
\norm{v}^2=\sum_{i=1}^n\abs{\langle v,u_i\rangle}^2.
\end{equation*}
\end{Theorem}
\begin{proof}
Write $v=\sum_i\alpha_iu_i$. Taking the inner product with $u_j$ gives $\alpha_j=\langle v,u_j\rangle$. Orthogonality and Pythagoras then give the norm identity.
\end{proof}

\begin{Definition}
For $S\subseteq V$, its \newterm{orthogonal complement} is
\begin{equation*}
S^\perp=\qty{\, v\in V \mid \langle v,s\rangle=0\text{ for every }s\in S \,}.
\end{equation*}
\end{Definition}

The set $S^\perp$ is a subspace and equals $(\operatorname{span}S)^\perp$. In a finite dimensional inner product space, the Gram--Schmidt process below gives
\begin{equation*}
V=W\oplus W^\perp,
\qquad
\dim W+\dim W^\perp=\dim V,
\qquad
(W^\perp)^\perp=W
\end{equation*}
for every subspace $W$.

\subsection{Gram--Schmidt Orthogonalisation}

\begin{Theorem}[Gram--Schmidt]\label{thm:GramSchmidt}
Let $(v_1,\ldots,v_k)$ be a linearly independent list in an inner product space. Define recursively
\begin{align*}
w_1&=v_1,& u_1&=\frac{w_1}{\norm{w_1}},\\
w_j&=v_j-\sum_{i=1}^{j-1}\langle v_j,u_i\rangle u_i,& u_j&=\frac{w_j}{\norm{w_j}}\qquad(2\le j\le k).
\end{align*}
Then $(u_1,\ldots,u_k)$ is orthonormal and
\begin{equation*}
\operatorname{span}\{u_1,\ldots,u_j\}=\operatorname{span}\{v_1,\ldots,v_j\}
\end{equation*}
for every $j$. Consequently every finite dimensional inner product space has an orthonormal basis.
\end{Theorem}
\begin{proof}
Suppose $u_1,\ldots,u_{j-1}$ are orthonormal and span the first $j-1$ vectors. For $r<j$,
\begin{equation*}
\langle w_j,u_r\rangle=\langle v_j,u_r\rangle-\sum_{i=1}^{j-1}\langle v_j,u_i\rangle\langle u_i,u_r\rangle=0.
\end{equation*}
Moreover, $w_j\ne0$, since otherwise $v_j$ would lie in the span of its predecessors. Hence $u_j$ is defined, has norm $1$, and is orthogonal to the preceding vectors. Since $v_j=w_j+\sum_{i<j}\langle v_j,u_i\rangle u_i$, adjoining $v_j$, $w_j$, or $u_j$ produces the same span. Induction proves the assertions. Apply the construction to any basis of a finite dimensional space.
\end{proof}

\begin{Example}
Apply Gram--Schmidt in $\mathbb R^3$ to
\begin{equation*}
v_1=\begin{bmatrix}1\\1\\0\end{bmatrix},\qquad v_2=\begin{bmatrix}1\\0\\1\end{bmatrix},\qquad v_3=\begin{bmatrix}0\\1\\1\end{bmatrix}.
\end{equation*}
The first two steps give
\begin{equation*}
u_1=\frac1{\sqrt2}\begin{bmatrix}1\\1\\0\end{bmatrix},\qquad
u_2=\frac1{\sqrt6}\begin{bmatrix}1\\-1\\2\end{bmatrix}.
\end{equation*}
The third gives $w_3=\frac23\begin{bmatrix}-1\\1\\1\end{bmatrix}$ and hence
\begin{equation*}
u_3=\frac1{\sqrt3}\begin{bmatrix}-1\\1\\1\end{bmatrix}.
\end{equation*}
Thus $(u_1,u_2,u_3)$ is an orthonormal basis with the same successive spans as $(v_1,v_2,v_3)$.
\end{Example}

\subsection{Orthogonal Projections and Unitary Matrices}

\begin{Theorem}\label{thm:OrthogonalProjection}
Let $(u_1,\ldots,u_k)$ be an orthonormal basis of a subspace $W$. The map
\begin{equation*}
P_W(v)=\sum_{i=1}^k\langle v,u_i\rangle u_i
\end{equation*}
is the projection onto $W$ along $W^\perp$. Thus $P_W^2=P_W$, $\im P_W=W$, $\ker P_W=W^\perp$, and $v-P_W(v)\perp W$.
\end{Theorem}
\begin{proof}
The displayed map is linear and takes values in $W$. For every $j$,
\begin{equation*}
\langle v-P_W(v),u_j\rangle=\langle v,u_j\rangle-\sum_i\langle v,u_i\rangle\langle u_i,u_j\rangle=0.
\end{equation*}
Hence $v=P_W(v)+(v-P_W(v))$ is the decomposition in $W\oplus W^\perp$, from which the remaining assertions follow.
\end{proof}

For a non-zero $v$, the projection of $u$ onto $\operatorname{span}\{v\}$ is
\begin{equation*}
\frac{\langle u,v\rangle}{\norm{v}^2}v,
\end{equation*}
and $u$ minus this projection is orthogonal to $v$.

\begin{Definition}
A square complex matrix $Q$ is \newterm{unitary} if $Q^*Q=I$, equivalently $Q^{-1}=Q^*$. A real unitary matrix is called \newterm{orthogonal}, so $Q^{-1}=Q^{\mathsf T}$.
\end{Definition}

\begin{Theorem}
For a square matrix $Q$, the following are equivalent:
\begin{enumerate}[label=(\roman*)]
\item $Q$ is unitary;
\item the columns of $Q$ form an orthonormal basis;
\item the rows of $Q$ form an orthonormal basis;
\item $\langle Qu,Qv\rangle=\langle u,v\rangle$ for all $u,v$;
\item $\norm{Qv}=\norm{v}$ for all $v$.
\end{enumerate}
\end{Theorem}
\begin{proof}
The entries of $Q^*Q$ are the inner products of columns, so (i) and (ii) are equivalent. Since a left inverse of a square matrix is also its inverse, $Q^*Q=I$ is equivalent to $QQ^*=I$, whose entries are the inner products of rows; this gives (iii). Moreover,
\begin{equation*}
\langle Qu,Qv\rangle=(Qv)^*Qu=v^*Q^*Qu,
\end{equation*}
proving the equivalence of (i) and (iv), and (iv) implies (v). Finally, (v) gives $v^*(Q^*Q-I)v=0$ for all $v$. The matrix $H=Q^*Q-I$ is Hermitian. Taking $v=e_i$, $e_i+e_j$, and $e_i+ie_j$ shows successively that every diagonal entry, every real part of an off-diagonal entry, and every imaginary part is zero. Thus $H=0$, proving (i).
\end{proof}

\subsection{Functionals on Inner Product Spaces}

For each $w\in V$, the map $f_w(v)=\langle v,w\rangle$ is a linear functional. In finite dimensions every functional arises uniquely this way.

\begin{Theorem}[Finite-Dimensional Riesz Representation]\label{thm:RieszFinite}
Let $V$ be a finite dimensional inner product space. For every $f\in V^*$, there is a unique $w\in V$ such that
\begin{equation*}
f(v)=\langle v,w\rangle\qquad(v\in V).
\end{equation*}
\end{Theorem}
\begin{proof}
Choose an orthonormal basis $(u_1,\ldots,u_n)$ and put
\begin{equation*}
w=\sum_{i=1}^n\overline{f(u_i)}u_i.
\end{equation*}
If $v=\sum_i\langle v,u_i\rangle u_i$, then
\begin{equation*}
\langle v,w\rangle=\sum_i\langle v,u_i\rangle f(u_i)=f(v).
\end{equation*}
If both $w$ and $w'$ represent $f$, then $\langle v,w-w'\rangle=0$ for every $v$. Taking $v=w-w'$ gives $w=w'$.
\end{proof}

\begin{Remark*}
The Riesz map $w\mapsto f_w$ is conjugate linear over $\mathbb C$, because $f_{\alpha w}=\overline\alpha f_w$ under our convention that the first argument is linear. This is different from the basis-dependent linear isomorphism $V\cong V^*$ in \cref{thm:DualBasis}.
\end{Remark*}

